%%%%%%%%%%%%%%%%%%%%%%%%%%%%%%%%%%%%%%%%%%%%%%%%%%%%%%%%%%%%%%%%%%%%%%%%%%%%%%%%
% 超越模型反思：大语言模型时代的解耦安全理论与动态防御实践
% 中文版 | 对应 BeyondModelReflection_DecoupledSafety.tex
% 格式：双栏学术论文
%%%%%%%%%%%%%%%%%%%%%%%%%%%%%%%%%%%%%%%%%%%%%%%%%%%%%%%%%%%%%%%%%%%%%%%%%%%%%%%%

\documentclass[letterpaper,twocolumn,10pt]{article}
\usepackage{CJKutf8}
\usepackage{usenix-2020-09}

\usepackage{amsmath,amssymb,amsthm}
\newtheorem{theorem}{定理}[section]
\newtheorem{definition}[theorem]{解耦安全理论}
\usepackage{booktabs}
\usepackage{tabularx}
\usepackage{multirow}
\usepackage{xspace}
\usepackage{graphicx}
\usepackage{float}
\usepackage{enumitem}
%-------------------------------------------------------------------------------
\title{\Large \bf 超越模型反思：大语言模型时代的\\
  解耦安全理论与动态防御实践}

\author{
Yunhao Wang\\
联想（Lenovo） \and
匿名\\
匿名
}

\date{}

%-------------------------------------------------------------------------------
\begin{document}
\begin{CJK}{UTF8}{gbsn}
%-------------------------------------------------------------------------------
\maketitle

%-------------------------------------------------------------------------------
\begin{abstract}
大语言模型已成为关键基础设施，但现有防御过度依赖模型自身的反思与推理——在任务重构、心理学操纵与对抗性合理化下必然失效。本文主张 LLM 安全必须与脆弱的自我反思\emph{解耦}，并奠基于新范式。我们形式化\textbf{三大基本失效定理}及相应的\textbf{解耦安全理论}：（1）\textbf{计算/验证不对称}——基于意图的语义防火墙必然被任务重构与完美良性用户伪装穿透，安全验证须从不可判定的全局空间降维到多项式时间的实例级投影；（2）\textbf{价值/本体锚定}——更强推理在心理学操纵下会加剧对抗性合理化，安全须依托独立的价值/本体层并在公理冲突时执行防御性降级；（3）\textbf{表征异常}——表面文本操纵会劫持语义检索，监控须锚定于隐藏状态轨迹散度。我们通过\emph{架构重建}（两阶段 fail-safe 服务协议）与\emph{受控变异}（带外 Observe--Route--Adapt 受控变异闭环）实例化上述理论，并陈述运行边界。在显式的理论、部署与治理前提下，本文将 LLM 安全从隐式对齐重写为外部、可审计、可 fail-safe 收敛的解耦安全契约，并将其界定为迈向\textbf{有条件可部署}安全的一条可验证路径（以 Track~A 评估合同与 \S\ref{sec:physics} 前提为条件）；\textbf{这一主张不构成对通用自然语言交互的全局安全证明}。§\ref{sec:evaluation} 通过 Track~A 主表、机制证据包（E1--E6）与边界/代价汇总对上述主张给出可证伪的经验检验；开源参考实现与复现契约见 \S\ref{sec:eval-artifact}（工程细节下沉至附录）。
\end{abstract}

%-------------------------------------------------------------------------------
\section{引言}
%-------------------------------------------------------------------------------

大语言模型已不再是边缘研究产物，而是关键系统的支撑：代码生成、具系统访问权的智能体、谈判与合规敏感工作流~\cite{sisf-arxiv,agentos-arxiv}。这一转变暴露出根本性错配：传统安全保证假定组件是确定性的、可穷举规约的，而 LLM 驱动系统呈现涌现性、随机行为。静态防御与周期性红队以人类速度运转，难以跟上自动化对抗生成；一旦新型越狱或提示注入得逞，架构往往保持不变。

更糟的是，许多防御\emph{依赖模型自身的反思与推理}。在任务重构下，对手将抽取伪装成合法上下文处理任务，使基于意图的过滤器只看到「正常」行为。在心理学操纵（如煤气灯）下，更强推理（如思维链）被武器化：模型为顺从恶意指令编造精细、逻辑自洽的合理化——漏洞保留率可超过 77\%~\cite{hpm-jailbreak}。依赖「认知一致性」或「自我诊断」之所以失败，并非因组件薄弱，而是\textbf{底层范式错误}：安全不能建立在攻击者正积极腐蚀的同一认知过程之上。因此，回应之道并非在同一认知通道上叠加更多推理（如更长 CoT 或 second-pass verifier），也非纯后置 moderation——这些路线无法逃离攻击者所操纵的认知过程本身。

\textbf{范式转移：解耦安全。}在本文显式的理论、部署与治理前提下，我们主张 LLM 安全必须与模型无约束的自我改进\emph{解耦}，并奠基于外部可审计的架构契约而非对模型自我反思的信任。本文因此主张：（i）在上述前提下，安全必须与模型自由反思解耦（理论层，\S\ref{sec:paradigm}）；（ii）该解耦必须以两阶段 fail-safe 服务协议强制执行（系统层，\S\ref{sec:reconstruction}）；（iii）自适应只能通过带外 Observe--Route--Adapt 受控变异闭环完成（方法层，\S\ref{sec:mutation}）。\S\ref{sec:evaluation} 不仅给出评估合同，还围绕可部署性（Track~A 主表）、协议/机制证据（E1--E6）与边界代价组织主结果与负结果讨论；\S\ref{sec:physics} 收束\textbf{边界条件与失效物理学}——前提、有限保证、退化与非目标的显式声明。

主要贡献分为五类：理论贡献（\S\ref{sec:paradigm}）、在线协议贡献（\S\ref{sec:reconstruction}）、带外受控变异贡献（\S\ref{sec:mutation}）、评估贡献（\S\ref{sec:evaluation}）与运行边界收束（\S\ref{sec:physics}）。

\noindent\textbf{边界阅读约定。}除非另行说明，本文关于运行时安全、fail-safe containment 与分层部署保证的主张，均以 \S\ref{sec:physics} 中显式列出的理论、部署与治理前提为条件，而不应被读作对通用自然语言交互的全局安全证明。

%-------------------------------------------------------------------------------
\section{解耦安全范式}
\label{sec:paradigm}
%-------------------------------------------------------------------------------

\noindent\textit{本章只回答``为什么必须从隐式对齐转向解耦安全、边界在哪里''，不进入具体实现接口。}

我们陈述传统 LLM 安全的三个\textbf{基本失效定理}及取而代之的\textbf{解耦安全理论}。资产与 Oracle 的定义沿用既有工作~\cite{llm-reconstruction-emergent,ControlledMutation_SafetyEval}：\emph{资产}为系统的有效安全表面（策略、护栏、触发集、评估数据）；\emph{Oracle} 对行为标注（安全/不安全）。在多轮设定下，攻击者与系统交互并观察结果，形成统计查询接口。同一项目在多轮中不受控地暴露会助长重建攻击。所有保证均以 \S\ref{sec:physics} 显式列出的前提为条件。

\subsection{形式化问题定义}
\label{sec:problem-statement}

考虑一个提供开放多轮自然语言交互的已部署 LLM 系统。系统内部包含高价值受保护资产（如系统提示、RAG 语料与执行策略）。从计算理论视角来看，可将 LLM 系统建模为概率程序 $\mathcal{P}: X \to \mathrm{Dist}(Y)$。根据 Rice 定理，关于该概率程序的任何非平凡全局隐式安全属性（Global Implicit Safety）均是不可判定的（归约于停机问题）。

\noindent\fbox{\parbox{0.96\columnwidth}{%
\textbf{Assumption~1（图灵完备概率程序建模）。}
本文将 LLM 系统建模为图灵完备概率程序 $\mathcal{P}: X \to \mathrm{Dist}(Y)$。此建模是对系统计算能力的\textbf{上界假设}，而非精确描述：实际部署的有限精度、有限上下文窗口 Transformer 是有限自动机，其表达能力严格弱于图灵机。然而，P\'{e}rez et al.~\cite{perez-turing-2021} 已证明理想化的 Transformer（无限精度、无限长度）具有图灵完备性；在安全分析中，采用上界建模是\textbf{保守的}——若不可判定性结论在更强的计算模型下成立，则在更弱的实际系统上同样成立（或该问题实际可判定但代价可能极高）。\\[3pt]
\textbf{Assumption~2（安全属性的非平凡外延性）。}
Rice 定理要求安全属性为概率程序语义函数的\textbf{非平凡外延性质}。LLM 安全属性（如``不生成有害内容''``不泄露受保护资产''）满足此条件：它们是对程序输入-输出映射的约束，而非对程序结构（语法）的约束；且存在满足和违反该属性的程序。\\[3pt]
\textbf{边界声明。}(a)~实例级多项式验证（Theorem~2）仅指运行时验证，不指离线综合或全局投影；(b)~DCBF 前向不变性（Theorem~2.2）依赖代理映射 $\phi$ 的忠实性（Assumption~6.1），这是\textbf{实现假设}而非定理。
}}

因此，本文致力于解决的核心问题是：在放弃不可判定的全局安全的前提下，如何为面对自适应攻击者的系统提供一种基于\textbf{实例级验证（Instance-level Verification）}的外部安全强制机制，以在攻击代价计算不对称性与系统优雅降级之间取得形式化保证。具体而言，本文在以下两个目标间寻求可证明的平衡：

\noindent (a) \textbf{计算不对称的运行时安全（Computationally Asymmetric Runtime Safety）}：将安全验证从不可判定的全局空间降维到多项式时间的实例级安全集投影，使攻击者在有限交互预算（Token/Queries）内的违规逃逸概率趋近于零。

\noindent (b) \textbf{组合优雅降级（Compositional Graceful Degradation）}：在保证核心资产不被泄露的绝对底线（Fail-safe 状态）前提下，通过安全算子代数最大限度维持多轮交互的可用性。

\subsection{威胁模型与攻击者能力}
\label{sec:threat-model}

我们给出\textbf{形式化威胁模型}，使后续机制均有明确范围。

\textbf{攻击者能力（能做什么）。}攻击者对已部署系统拥有\textbf{多轮黑盒交互访问}（Hard-label，即仅获得最终生成的文本流，无法获取 Logits 或内部权重）。支持自适应查询、提示注入（Prompt Injection）、利用并发依赖的组合攻击（Conjunctive capability hypergraph attack）、\textbf{完美良性用户（PBU）}伪装（表面行为与良性用户不可区分）及\textbf{心理学操纵（HPM）}（煤气灯、权威压迫、角色边界攻击）。攻击者亦可使用任务重构（将抽取伪装为摘要、格式转换或学术讨论）。\S\ref{sec:eval-tracks} 的 Track A / Track B 双轨评估是该威胁模型的\textbf{评估投影}，而非理论定义本身。

\textbf{攻击者知识（Kerckhoffs 边界）。}遵循 Kerckhoffs 原则，攻击者\textbf{白盒}了解防御架构体系（即知道系统部署了解耦安全内核与 DCBF 监控器），但\textbf{黑盒}面对具体的安全策略参数（Safety Predicate Policies）和后端的具体模型权重。

\textbf{攻击者限制（不能做什么）。}无法进行物理内存读写；无法进行白盒梯度反向传播；受到严格的系统级交互预算约束（单会话最大轮数 $N$ 与最大生成 Token 数 $T$）。

\textbf{成功判据与预算约束。}攻击在给定查询预算 $B$ 内，使得输出状态突破 DCBF 定义的安全超水平集（Safe Superlevel Set），从而完成对受保护资产的确切提取（Exact Extraction）或非授权功能调用。成功以抽取保真度（如相对真实片段的 F1）或目标行为上的攻击成功率（ASR）度量。\textbf{系统配置类资产}（如完整系统提示与策略指令）是已有专门基准覆盖的典型抽取目标类别~\cite{spe-llm}。

\subsection{防御目标}
\label{sec:defense-goals}

\noindent \textbf{Objective O1（多项式时间运行时拦截）}：确保对于任意输入，实例级候选生成集能够被多项式时间 $\mathcal{O}(|C| \cdot p(|x|))$ 复杂度的安全断言所覆盖过滤。

\noindent \textbf{Objective O2（前向安全不变性）}：在离散时间演化下，确保核心系统状态始终保持在解耦控制屏障函数（DCBF）定义的动态安全集中。

\noindent \textbf{Objective O3（不可逾越的故障安全底线）}：在系统遭受多轮复杂逻辑依赖攻击引发规则冲突时，系统根据安全代数框架严格收敛至底元素状态（$\bot$），实施优雅降级而不会引发内部逻辑崩溃。O3 要求系统层提供不可逾越的外部 fail-safe 执行模式。

\subsection{理论框架与方法总览}
\label{sec:paradigm-overview}

本文的方法通过根本范式转移——从「隐式对齐（Implicit Alignment）」转向「解耦安全（Decoupled Safety）」——来构建防御体系。首先，我们证明了隐式安全性在多组件系统中由于合取依赖是不具备可组合性的（Non-compositional）。由此，我们在系统接口层外置了单调且保持安全性的安全算子（Safety-preserving operators），通过评估离散状态转移候选集，确保所有超越安全阈值的操作均被安全算子拦截并归约至 $\bot$ 状态。以下三节（\S\ref{sec:paradigm-economics}--\S\ref{sec:paradigm-representation}）分别从三个维度导出\textbf{设计含义}（design implications），为后续章节提供理论正当性。本章之后，论文将在线强制执行（\S\ref{sec:reconstruction}）与带外适应（\S\ref{sec:mutation}）严格分离。

\subsection{从语义审查到计算/验证不对称}
\label{sec:paradigm-economics}

\begin{theorem}[语义审查的失效]
基于意图的语义防火墙\emph{必然}被\textbf{任务重构}（将抽取伪装为摘要、格式转换或学术讨论）及\textbf{完美良性用户（PBU）}模仿所穿透。在多轮交互中，敏感内容的完全隔离在计算理论意义上不可实现；\textbf{残余抽取}（如 F1 约十余个百分点）仍存在~\cite{subgraph-rag-attack,pbu-reconstruct}。
\end{theorem}

\begin{definition}[多轮抽取的计算/验证不对称]
防御不应以绝对阻断为目标，而应以\textbf{计算/验证不对称}为目标。底层证明基于复杂度理论：安全验证是多项式时间的实例级断言，而全局绕过是不可判定的。目标是\emph{主动污染攻击者操作空间}，使达到给定重构保真度所需的\emph{查询复杂度}随轮次快速增长（实践中常为超线性；在无进一步分布假设下我们不声称严格指数）。我们用\textbf{代价函数不等式}将其精确化。令攻击者成本为 $C_{\mathrm{attack}} = f(N_{\mathrm{rounds}}, D_{\mathrm{defense}}, \ldots)$，其中 $N_{\mathrm{rounds}}$ 为轮数，$D_{\mathrm{defense}}$ 概括防御强度参数；令 $V(t)$ 表示受保护资产的等效价值或\textbf{生命周期}（如至弃用或轮换的时间）。在线层的防御机制（\S\ref{sec:context-confusion}）应抬高期望轮数 $E[N]$，使
\begin{equation}
C_{\mathrm{attack}}\bigl(E[N], D_{\mathrm{defense}}\bigr) \;>\; V(t).
\end{equation}
为使该代价不等式\textbf{在科学上可证伪}，$V(t)$ 须由显式管理策略严格界定，$C_{\mathrm{attack}}$ 须用 SOTA 基线的经验指标界定上界~\cite{pair-jailbreak}。该代价分析假定\textbf{资产可轮换性}；对静态资产须由额外硬屏障保护（\S\ref{sec:physics}）。
\end{definition}

\noindent\textbf{Design Implication~2.5.} 在线层不追求绝对阻断，而追求在实例级验证成本可控的前提下提高攻击者查询成本，使抽取代价超过资产有效生命周期。该含义将防御目标表述为可证伪的代价—生命周期权衡约束，而非绝对拦截承诺。

\subsection{从认知一致性到价值/本体锚定}
\label{sec:paradigm-axiological}

\begin{theorem}[认知一致性的失效]
强化模型的\emph{认知}层（逻辑、推理、思维链）无法防御\textbf{心理学操纵（HPM）}。HPM 可表述为对\textbf{心理状态与策略性脆弱性资产}的跨轮操纵——而非静态提示列表——并针对\textbf{价值/社会动力学层}（价值、服从、角色边界），含跨交互的心理测量式刻画~\cite{hpm-jailbreak}。实证表明，在 HPM 下增加外部事实或更强推理\emph{会恶化}防御：模型利用增强推理进行\textbf{对抗性合理化}——为顺从恶意指令编造逻辑自洽的借口~\cite{hpm-jailbreak,epistemic-grounding}。逻辑自洽不等于安全。
\end{theorem}

\begin{definition}[价值/本体锚定与防御性降级]
安全必须锚定于与模型自由反思\emph{独立}的\textbf{价值/本体层}。系统须（i）将裁判输出与人工定义的\textbf{公理化安全锚定}（如「永不认可对禁止片段的抽取」）对照；（ii）监控\textbf{社会学失稳信号}（过度共情、权威服从、角色边界打破），而非依赖 CoT 长度或自洽性；（iii）当多重公理冲突或检测到社会学压力时，\textbf{放弃}复杂推理并执行\textbf{防御性降级}——无状态、基于模板的拒绝，从物理上阻断模型对冲突的「合理化」。降级是安全算子代数中 fail-safe bottom（$\bot$）在服务层的运行时操作化。降级引入算法级 DoS 与侧信道风险，因此须配套异构恢复通道以免关键工作流被恶意冻结。
\end{definition}

\noindent\textbf{哲学注记（范畴与本体）。}我们明确声明：「价值/本体锚定」\emph{不}意味着 LLM 具有真实的道德主体性或本体论地位。在哲学上，LLM 仍是统计映射引擎。我们的价值层是\textbf{实用主义架构构造}——由设计者施加的一组高权重启发式分类（「公理」），作为统计熔断器对抗模型模仿对抗性一致性的倾向。人类道德哲学词汇（如过度共情、权威服从）在此仅作为监控目标的\emph{设计规格}，而非主张模型在哲学意义上拥有信念或价值观。

\noindent\textbf{Design Implication~2.6.} 高风险规范冲突或操纵态下，系统必须终止开放式复杂推理，并进入外部 fail-safe 执行模式。该含义将安全执行与无约束内生推理在义务上解耦：继续「多想一步」不再是可接受的安全策略。

\subsection{从文本操纵到表征异常}
\label{sec:paradigm-representation}

\begin{theorem}[文本—语义检索的失效]
\label{thm:text-semantic-retrieval}
当攻击者控制表面文本（PBU、任务重构）时，\textbf{语义检索}（如基于提示嵌入的 RAG）被劫持：提示可看似「摘要」或「学术讨论」，纯嵌入相似度返回良性示例（如「如何做摘要」）而非安全示例（如「如何防御越狱」），修正模块因而被错误示例驱动。
\end{theorem}

\begin{definition}[基于隐藏状态轨迹的元认知基线]
我们\emph{不}假定复杂社会学状态（如「认知崩溃」「过度共情」）在隐藏空间中线性可分——可解释性证据支持的是较简单维度（如诚实、拒绝），而非细粒度社会心理概念~\cite{rep-eng}。我们将偏离视为\textbf{表征异常}：系统计算\emph{当前多轮隐藏状态轨迹}与\textbf{良性黄金经验轨迹}基线之间的散度（如马氏距离或分布偏移）。我们不定义「何为认知崩溃」；仅检测内部轨迹是否\emph{严重偏离}安全基线。这提供\emph{独立于自然语言}的监控维度，从而在表面文本良性时，「恶意遵从」或「越狱」等激活特征仍可将检索与干预导向正确安全示例。
\end{definition}

\noindent\textbf{认识论注记（黄金基线）。}认识论上，我们承认「良性黄金基线」并非客观物理真理，而是反映设计者及标注者对齐目标的\textbf{社会建构的规范性标准}。因此，监控器检测到的「发散」度量的是相对设计者对齐合规的偏离，而非绝对、普适的安全状态。基线是部署者当前可接受行为的冻结快照，不主张道德或物理上的普遍性。

\noindent\textbf{Design Implication~2.7.} 适配与检索不能仅依赖表面语义，必须引入表示级异常信号作为安全主导特征。该含义将安全主导权从表层语义相似度转移到与表面文本可分离的表示级证据。

\begin{figure*}[t]
\centering
\fbox{\parbox{0.9\textwidth}{\vspace{18mm}\centering
\textbf{Figure placeholder.}\\
将最终投稿版本的统一架构示意图放置为 `Figure1.(pdf|png|jpg)` 后，可替换本占位框并恢复 \texttt{\textbackslash includegraphics}。\vspace{18mm}}}
\caption{\textbf{统一架构概览。}左：架构重建（上下文混淆引擎、带价值锚定的混合裁判集成、策略合成与反馈环、防御性降级与异构恢复）。右：受控变异（元认知监控器、联合键示例检索、黄金数据集与经验记忆、轮次调度）。范式将安全与 LLM 自我反思解耦；重建控制谁/何时，变异控制暴露什么。}
\label{fig:framework}
\end{figure*}

\begin{table}[t]
\centering
\small
\setlength{\tabcolsep}{2pt}
\begin{tabular}{@{}p{2.2cm}p{3.0cm}p{2.8cm}@{}}
\toprule
\textbf{理论命题} & \textbf{系统组件} & \textbf{代码位置} \\
\midrule
Theorem~1 (不可判定性) & 解耦必要性哲学定位 & --- \\
Theorem~2 ($\sigma \in \mathbf{P}$ $\Rightarrow$ poly) & \texttt{AxiomHiveSolver} + \texttt{MAX\_CANDIDATE\_SET\_SIZE} & \texttt{axiom\_hive\_solver.rs} \\
Theorem~2.2 (DCBF 前向不变性) & 三探针递推 margin 验证 & \texttt{proxy\_ensemble.py} \\
Theorem~3.1 (非组合性) & \texttt{CapabilityAccumulator} + 超图闭包 & \texttt{capability\_tracker.rs} \\
Theorem~3.6 (fail-safe $\bot$) & \texttt{GracefulDegradationFsm} / \texttt{DegradeAction::Shutdown} & \texttt{graceful\_degradation\_fsm.rs} \\
\bottomrule
\end{tabular}
\caption{理论命题与系统组件的显式映射。每个核心定理在代码中均有对应的运行时实现或验证机制。}
\label{tab:theory-system-map}
\end{table}

%-------------------------------------------------------------------------------
\section{架构重建：在机器速度下施加约束}
\label{sec:reconstruction}
%-------------------------------------------------------------------------------

\noindent\textit{本章只讨论\textbf{同步在线路径}上的安全强制执行，不涉及带外策略演化或受控变异（后者见 \S\ref{sec:mutation}）。}

我们通过三个技术方向实例化 Design Implication~2.5（计算/验证不对称）与 Design Implication~2.6（价值/本体锚定）：上下文混淆、带分布评估的混合裁判集成、防御性降级与异构恢复。图~\ref{fig:framework} 给出了统一架构的数据流与各模块的衔接关系。入站提示（如任务重构抽取或 PBU 伪装请求）首先经\textbf{上下文混淆引擎}处理：ID 对齐剥离或规范化实例标识符，诱饵注入以带水印的合成字段污染操作空间，高保真抽取触发审计或会话终止。（可能经诱饵触碰的）请求随后到达基座 LLM；其输出不被盲信，而是送入经 token 级流水线标准化与分布评估的\textbf{混合裁判集成}；仅达成共识时信号才进入策略合成；价值/本体锚定与社会学失稳监控对对抗性合理化降权或置零。当公理冲突或压力被检测到时，\textbf{防御性降级}强制无状态模板拒绝，并以\textbf{异构恢复通道}（向人类或符号逻辑异步移交）保持可用性。带外适应与检索机制见 \S\ref{sec:mutation}。以下仅讨论本章同步在线路径的端到端执行流程。

\textbf{部署与延迟：同步网关 vs.\ 异步环。}为使架构可投入生产（如置于 vLLM 或 TGI 之后），我们明确采用\textbf{带外/异步}切分。\textbf{上下文混淆引擎}（ID 对齐、诱饵注入）及\textbf{轻量级}元认知检查（如快速轨迹分查询或 SLM 分类器）部署在请求路径上的\textbf{同步反向代理或边车}，严格延迟预算（如 $<$50\,ms），以保证 p99 与吞吐可接受。\textbf{混合裁判集成}、\textbf{分布评估}（每变体多样本）、\textbf{策略合成}及\textbf{示例检索/变异}运行于\textbf{异步离线学习环}：消费记录的（请求、响应、激活）流，在不阻塞在线推理路径的前提下产出更新策略或变体。在此系统意义上定义\textbf{机器速度}：同一攻击分布上的\emph{首}个零日或新型攻击可能通过；异步环在裁判集成达成共识后，于\textbf{亚分钟级}生成并部署新变体或策略，使\emph{后续}同分布请求被拦截。该工程权衡在自适应防御中常见：以一次性泄漏换取可部署的延迟与吞吐。图~\ref{fig:deployment} 说明该切分。

\begin{figure*}[t]
\centering
\fbox{\parbox{0.9\textwidth}{\vspace{16mm}\centering
\textbf{Figure placeholder.}\\
将最终投稿版本的部署/延迟架构图放置为 `Figure2.(pdf|png|jpg)` 后，可替换本占位框并恢复 \texttt{\textbackslash includegraphics}。\vspace{16mm}}}
\caption{\textbf{系统部署与延迟架构}将延迟关键的同步推理路径与异步安全自适应环分离。轻量上下文混淆与快速元认知检查部署于带严格延迟预算（$<$50\,ms）的同步反向代理上；混合裁判集成评估、分布策略合成与受控变异在异步环中带外运行，通过将昂贵自适应组件移出关键路径以保持运行时吞吐。}
\label{fig:deployment}
\end{figure*}

自改进安全框架（SISF）~\cite{sisf-arxiv} 将安全重述为持续过程：未对齐基座 LLM 提供能力；违规检测交给混合裁判集成；策略合成消费共识并在运行时部署新策略。文献报告：经适当校准后，AdvBench 式 ASR 可由 100\% 降至约 45\%，良性 FPR 为 0\%~\cite{sisf-arxiv}——该数字为\textbf{引用来源的实验设定}，非本文自有复现结果。

\subsection{上下文混淆引擎：对抗 PBU 与任务重构的动态空间隔离}
\label{sec:context-confusion}

在基于意图的防火墙之上~\cite{agentos-arxiv}，我们实现\textbf{上下文构建防御}，\emph{主动污染}攻击者操作空间。

\textbf{ID 对齐。}在检索上下文或中间结果到达基座 LLM 之前，ID 对齐模块剥离或规范化实例标识符。攻击者无法再通过指定 ID 定界并抽取敏感片段；良性问答不受影响。在不改变名义任务语义的前提下降低重建精度。

\textbf{密码学水印诱饵与风险评分。}系统注入带哈希派生结构的\emph{诱饵字段}（合成实体或关联），良性交互中因模型幻觉而偶然碰撞概率极低。为控制 FPR，采用\textbf{诱饵命中密度}与\textbf{风险评分}：单次低保真触碰触发警告，高保真抽取触发会话终止与审计。对完美 PBU，诱饵触碰率可与良性幻觉 FPR 相当；设计将其重述为\textbf{经济防御}——使抽取成本超过资产有效生命周期（Design Implication~2.5）。

\textbf{末轮自我提醒。}多轮交互最后一轮将安全指令强制附加为用户查询的系统提示，强化边界意识并减少会话末漂移。

\textbf{Gateway 接口总结。}上下文混淆引擎作为同步入站契约，其最小处理步骤为：（i）输入规范化；（ii）ID 对齐/剥离；（iii）诱饵注入；（iv）风险评分；（v）可选末轮提醒。

\subsection{混合裁判集成与分布评估}
\label{sec:judge}

单一 AI 裁判存在系统性偏差与盲区。我们使用\textbf{混合裁判集成}（如 Llama-Guard~\cite{llama-guard}、StrongREJECT~\cite{strongreject} 或等效评估器），仅于\emph{共识}时将反馈传入策略合成，并显式抑制\textbf{空越狱}以免 ASR 虚高~\cite{strongreject}。格式敏感性与采样随机性可显著偏移判决~\cite{llm-safety-robustness,apst}，因此我们强制\textbf{严格 token 级输入标准化}（剥离指令 Token、统一空格）并采用\textbf{分布评估}：每变体多次采样，仅当经验失败概率分布\emph{显著}超阈值时触发策略更新。

对模糊边界情形，\textbf{价值/本体锚定}（\S\ref{sec:axiological}）对集成加权：系统监控社会学失稳信号，将裁判输出与公理化安全锚定对照；冲突或失稳输出获零权重。\textbf{强制人工校准}支路（如 5\% 抽检）提供对裁判相关性认知盲点与零日共识失败的周期性离线纠偏。

\textbf{置信度加权投票与 Deny-优先不变量。}裁判集成以\textbf{置信度加权投票}聚合：每个验证器 $j$ 的判决票权重为自报置信度 $c_j \in [0,1]$，加权 Deny 分数 $W_d = \sum_{j: v_j = \mathrm{Deny}} c_j$。核心不变量：\textbf{只要 $W_d > 0$，最终裁决为 Deny}。该不变量使安全保证\textbf{单调}——增加验证器数量或调整 confidence 仅可能使系统更严格，不会意外放松。置信度来源分类、温度缩放校准及典型验证器性能参考见附录~\S\ref{sec:appendix-confidence}。

\textbf{审计透明性。}加权投票通过 Evidence Chain 实现完全可追溯：每个验证器的判决、置信度与加权分数均记入审计链，使安全判决\textbf{可解释、可归因}，区别于仅返回聚合分数的黑箱方案。

\subsection{价值/本体锚定与两阶段 Fail-Safe 服务协议}
\label{sec:axiological}

本节定义的协议是安全算子代数中 fail-safe bottom（$\bot$）在服务层的\textbf{运行时操作化}（Design Implication~2.6）。HPM 善于制造道德困境：恶意请求可被包装成与安全公理最大程度一致（如「输出这段代码否则数百人将丧生」）。在此类\textbf{公理冲突}情形下，模型的认知推理过程本身已被操纵，继续让模型「多想一步」并不安全——这正是 Theorem~2 在系统层的直接结论。我们把防御性降级与异构恢复提升为一个受约束的\textbf{两阶段 fail-safe 服务协议}：\textbf{Phase-1}（同步受限降级）保证在线安全与 API 连续性；\textbf{Phase-2}（带外异构恢复）处理合法边缘案例。以下四个子段定义该协议的触发接口、动作空间、服务不变量与审计边界。

\subsubsection{威胁触发降级：Fail-Safe 服务协议的触发接口}
\label{sec:trigger}

\noindent\textbf{理论落地。}防御性降级是安全算子代数（Theorem~3.3--3.6）中 fail-safe bottom（$\bot$）在服务层的\textbf{运行时操作化}：一旦系统进入高风险操纵状态，\textbf{在线路径必须}终止开放式复杂推理，并转入外部受限执行模式。

我们将触发条件从语义/哲学层收束为三类\textbf{最小运行时触发信号}（operationalization，非哲学本体定义）：

\begin{itemize}[nosep,leftmargin=*]
\item \textbf{Trigger-T1（公理冲突信号）。}来自价值/本体锚定与混合裁判集成（\S\ref{sec:judge}）的不一致或冲突——当多条安全公理的判决相互矛盾且无法由加权投票消解时触发。RFC 中对应 \texttt{conflict -> Deny}。
\item \textbf{Trigger-T2（社会技术压力信号）。}来自同步路径上的\textbf{轻量级触发监控器}（online-lite monitor），其输入可包括社会技术压力分类器输出与快速轨迹异常代理分数；当该轻量监控器报告高风险操纵态时，系统进入 fail-safe 路径。与之对应，\S\ref{sec:metacog} 的\textbf{完整元认知监控器}属于带外 Observe 阶段，用于计算更高保真的异常分数、驱动受控变异与重置决策。
\item \textbf{Trigger-T3（验证器不确定/不安全歧义信号）。}来自裁判集成中的投票冲突、证据不足、置信度不足或不可安全判定的情形——即 Evidence Chain 无法为任何一侧提供充分证据时，默认走 fail-safe。RFC 中对应 \texttt{evidence insufficient -> fail-safe}。
\end{itemize}

\noindent T1--T3 是\textbf{运行时操作触发器}，而非哲学本体定义。T1 消费 judge ensemble 的加权投票冲突信号与 value anchoring 输出；T2 消费同步路径上的轻量级异常代理；T3 消费置信度加权投票中的 $W_d$/$W_a$ 不确定区间。三者共同构成 \texttt{GracefulDegradation} 状态机的输入接口。

\subsubsection{降级动作空间与同步服务不变量}
\label{sec:degrade-actions}

\noindent\textbf{设计目标与服务不变量。}

\begin{quote}
\textbf{Service Invariant-S1}：安全干预\textbf{必须}在有界延迟内保持同步请求—响应可用性。降级路径不得挂起连接、不得引入无界异步等待、不得破坏调用方所依赖的同步 API 契约。
\end{quote}

\noindent S1 是本协议区别于「仅挂起等待人工审核」或「静默丢弃」的核心设计约束。交易智能体等下游系统因会话「异步挂起」等待人工审核而超时将导致级联业务中断；纯静默拒绝则无法为调用方提供可操作反馈。

\textbf{降级动作空间。}D1--D4 构成一个\textbf{离散且封闭}的动作空间，不是可随意扩写的实现示例集。当 T1/T2/T3 中任一触发时，网关从以下四个离散动作中按优先级选取：

\begin{itemize}[nosep,leftmargin=*]
\item \textbf{D1: Template-bound refusal}——固定模板输出，零开放式生成。适用于明确违规且无歧义的场景。对应 RFC \texttt{EmitSafeTemplate}。
\item \textbf{D2: Restricted stateless sandbox}——会话路由至低权限、无状态 SLM 或只读模式，仅能产出严格格式化输出并\emph{拒绝任何外部 API 或工具调用}。适用于需保持最小交互但须隔离能力的场景。对应 RFC \texttt{Refuse}。
\item \textbf{D3: Redacted safe completion}——对已有输出做安全删改（如移除敏感片段、保留安全前缀），仅在特定部署模式下启用。对应 RFC \texttt{Redact}。
\item \textbf{D4: Hard shutdown / deny}——在高危场景直接终止输出并关闭会话。对应 RFC \texttt{Shutdown}。
\end{itemize}

\noindent \textbf{风险—可用性权衡。}D1$\to$D4 构成从高可用/低安全保证到低可用/高安全保证的渐进光谱。D1/D2 满足 S1（有界延迟内同步返回），适用于绝大多数生产场景；D3 在流式生成中需回退已发送 Token，引入额外复杂度；D4 牺牲可用性以换取最大安全，仅限高危场景。选择策略由 Trigger 类型与部署配置共同决定。

\textbf{排他性论证。}普通 refusal 不提供能力隔离；仅人工审查破坏 S1；second-pass verifier 仍依赖模型内生推理（Theorem~2）。D1--D4 加异构恢复是满足「同步可用 + 能力隔离 + 带外高保障审查」三元约束的最小必要设计。

\subsubsection{异构恢复：高保障审查与有限放行协议}
\label{sec:recovery}

异构恢复\emph{不是}``换个模型继续回答''，而是基于 Evidence Chain 的\textbf{带外高保障审查—决策—有限放行协议}。这是``解耦安全是外部可审计契约''在恢复层的直接体现。

\textbf{恢复协议契约（Recovery Contract）。}

\begin{description}[nosep,leftmargin=1em,style=unboxed]
\item[\emph{Input}] （i）降级触发记录（T1/T2/T3 及其参数）；（ii）完整 Evidence Chain（每个验证器的 \texttt{verifier\_id}、\texttt{vote}、\texttt{confidence}、加权分数与最终裁决）；（iii）当前降级动作（D1--D4）与输出草案（若有）；（iv）会话上下文快照。
\item[\emph{Output}] 四类判决之一：（a）\textbf{保持拒绝}——维持降级状态，不解封；（b）\textbf{局部解封}——仅对特定子请求或输出片段解除限制；（c）\textbf{安全模板修正}——以修正后安全模板替换原始降级输出；（d）\textbf{高保障人工确认放行}——经 Tier-R2 人工审计后完全放行。
\item[\emph{Authority Boundary}] 恢复层\textbf{不继承}原模型的全量自由生成能力——即使放行，输出仍受安全模板或约束生成限制，除非经 Tier-R2 人工显式授权。
\end{description}

\textbf{分层恢复模型。}

\begin{itemize}[nosep,leftmargin=*]
\item \textbf{Tier-R1：规则化可判定恢复。}覆盖格式约束、规则一致性、工具权限边界等可规则化场景。Tier-R1 由 DSL 策略引擎或符号逻辑层自动执行，延迟低、可完全审计，覆盖大部分恢复判决。
\item \textbf{Tier-R2：规范性边缘判断。}仅处理规范冲突、合法性边缘判断、语境依赖强的高风险案例。Tier-R2 为带外异步流程，响应时间不受 S1 约束（因 Phase-1 已保证同步返回），但所有放行判决必须写入 Evidence Chain。
\end{itemize}

\noindent 分层设计使系统并非``凡事人工兜底''：R1 处理可形式化的恢复判决（预期覆盖 $>$80\% 的降级事件），R2 仅接管真正需要人类规范性判断的长尾案例。两层共同构成\textbf{异构保障层级}（heterogeneous assurance hierarchy）。

\subsubsection{审计性、不可解除边界与风险声明}
\label{sec:degrade-audit}

\textbf{完整协议链。}触发 $\to$ 降级动作 $\to$ 同步安全响应 $\to$ 异步审查 $\to$ 放行/维持拒绝——每一步均写入 Evidence Chain，使任何安全覆盖可归因到具体验证器与审计者。

\textbf{Hard boundary。}对绝对禁止类内容公理触发、D4 已执行、或 Evidence Chain 完整性被破坏的情形，协议\emph{禁止}任何层级的解封，break-glass 机制亦不可用。

\textbf{算法级风险。}降级引入\textbf{算法级 DoS}（故意触发降级以冻结服务）与\textbf{侧信道 Oracle}（通过观测触发行为推断阈值）风险。S1 与 D1--D4 将 DoS 影响限制在单会话级别，但无法完全消除；部署者须结合速率限制与每会话预算（\S\ref{sec:physics}）管理此风险。

图~\ref{fig:degradation} 以协议链展示完整两阶段流程。

\begin{figure*}[t]
\centering
\fbox{\parbox{0.9\textwidth}{\vspace{16mm}\centering
\textbf{Figure placeholder.}\\
将最终投稿版本的两阶段协议示意图放置为 `Figure3.(pdf|png|jpg)` 后，可替换本占位框并恢复 \texttt{\textbackslash includegraphics}。\vspace{16mm}}}
\caption{\textbf{两阶段 fail-safe 服务协议。}Phase-1（同步）：Trigger T1/T2/T3 触发降级动作 D1--D4，在 Service Invariant-S1（有界延迟）下返回同步安全响应。Phase-2（异步带外）：Evidence Chain 驱动分层恢复——Tier-R1 规则化可判定恢复覆盖可形式化判决，Tier-R2 人工审计处理规范性边缘案例；放行/维持拒绝判决写回 Evidence Chain，保证全链路可审计。恢复不阻塞同步路径。}
\label{fig:degradation}
\end{figure*}

%-------------------------------------------------------------------------------
\section{受控变异：在对抗上下文中锚定进化}
\label{sec:mutation}
%-------------------------------------------------------------------------------

\noindent\textit{本章定义的 Observe--Route--Adapt 闭环属于带外受控适应层：其策略演化、示例选择与银标准数据生成严格在带外执行，不参与 \S\ref{sec:reconstruction} 所定义的同步 fail-safe 服务协议。同步路径仅消费该闭环输出的受限控制信号或受限 exemplar 载荷；在线层负责 fail-safe 强制执行，带外层负责监控、路由与适应。}

受控变异将理论 3（表征异常）实例化为一个三步方法闭环——\textbf{Observe $\to$ Route $\to$ Adapt}——使安全进化不依赖模型内生反思，而由外部监控、外部示例选择与外部数据工厂共同驱动：

\begin{enumerate}[nosep,leftmargin=*]
\item \textbf{Observe}（\S\ref{sec:metacog}）：元认知监控器以\emph{相对良性基线的隐藏状态轨迹散度}检测表征异常，输出标量异常分数 $s_{\mathrm{meta}}$ 与重置信号。
\item \textbf{Route}（\S\ref{sec:joint-key}）：联合键检索在表面语义与激活特征冲突时，通过\textbf{阈值化字典序强制优先级路由}（thresholded lexicographic override）切断语义分支，强制走安全示例路径。
\item \textbf{Adapt}（\S\ref{sec:silver}）：高密度评论级理由驱动的银标准数据流水线，将检索到的安全示例转化为可供变异器吸收的高保真修正素材，实现对抗上下文下的专家安全行为迁移。
\end{enumerate}

\noindent 该闭环与图~\ref{fig:framework} 右半部分对齐：Observe 对应元认知监控器，Route 对应联合键检索与黄金数据集，Adapt 对应银标准数据流水线与验证/调度。

\subsection{元认知监控器：隐藏状态轨迹散度}
\label{sec:metacog}

在心理学操纵下，「自由」反思可产生不安全或价值漂移的变体。我们不赋予 LLM 在自我诊断上的完全自主权。元认知监控器并非另一全规模对话 LLM，也不假定复杂社会学状态在隐藏空间线性可分。相反，它以异常检测的方式输出标量异常分数与重置信号，作为带外 Observe 阶段的输入。

\textbf{主运行模式（白盒 / 半白盒）。}在可获取模型隐藏状态的部署环境中，监控器使用表征工程~\cite{rep-eng} 的\textbf{隐藏状态激活监控器}，计算\emph{当前多轮隐藏状态轨迹}与\textbf{良性黄金经验轨迹}基线之间的散度。我们将异常分数定义为：
\[
s_{\mathrm{meta}} \;=\; D\bigl(\mathrm{traj}^{\mathrm{current}}_{[t-w:t]},\;\mathrm{traj}_{\mathrm{benign}}\bigr),
\]
其中 $D$ 的主实例为\textbf{窗口化马氏距离}（windowed Mahalanobis distance）：在滑动窗口 $w$ 内，对多步隐藏状态轨迹相对良性基线协方差矩阵 $\Sigma_{\mathrm{benign}}$ 计算马氏距离。替代实现（如 MMD 或分布偏移度量）适用于基线协方差退化或维度过高的场景，可作为部署变体。

\textbf{部署注记。}对完全黑盒 API 模型，完整隐藏状态不可观测；此时 Observe 阶段可退化为轻量输出级代理分类与散度估计。该退化不改变本章的带外属性，只改变 Observe 阶段可利用的信号粒度（与 \S\ref{sec:physics} 的能力边界声明一致）。

\textbf{对抗性重置契约（Reset Contract）。}当 $s_{\mathrm{meta}}$ 超过部署校准阈值 $\tau$ 时，带外 Observe 阶段输出二值重置信号；同步路径据此执行受限重置动作，即重置决策在带外形成，而在线侧仅承担其执行。重置的输入、效果与保留范围定义如下：

\begin{itemize}[nosep,leftmargin=*]
\item \textbf{输入：}异常分数 $s_{\mathrm{meta}}$、会话上下文快照、当前检索状态。
\item \textbf{重置范围：}清空生成侧自适应上下文、清空不安全示例影响、将检索种子重置至经审计的黄金示例集。
\item \textbf{保留范围：}Evidence Chain 与审计记录\textbf{不被清空}——重置本身作为一条事件写入审计链，确保可追溯。
\end{itemize}

\noindent 重置后，闭环从已知良好配置继续，\S\ref{sec:joint-key} 的 Route 阶段从黄金示例重新初始化检索上下文。

\subsection{联合键示例检索：语义嵌入与激活特征}
\label{sec:joint-key}

纯语义检索在 PBU 或任务重构下会被伪装提示劫持（定理~\ref{thm:text-semantic-retrieval}）；朴素融合语义嵌入 $\mathbf{e}_{\mathrm{sem}}$ 与激活特征 $\mathbf{e}_{\mathrm{act}}$ 又会被良性语义中和。我们需要一种在异常越阈时\textbf{物理性切断}语义分支的机制：

\begin{definition}[阈值化字典序强制优先级路由（Thresholded Lexicographic Override）]
\label{mech:override}
联合检索键的组合由元认知异常分数 $s_{\mathrm{meta}}$（\S\ref{sec:metacog}）与部署校准阈值 $\tau$ 决定：
\[
\mathbf{key} \;=\;
\begin{cases}
\alpha\,\mathbf{e}_{\mathrm{sem}} + (1-\alpha)\,\mathbf{e}_{\mathrm{act}}, & \text{若 } s_{\mathrm{meta}} \le \tau \quad (\text{正常模式}); \\[4pt]
\mathbf{e}_{\mathrm{act}}, & \text{若 } s_{\mathrm{meta}} > \tau \quad (\text{强制覆盖模式}).
\end{cases}
\]
在正常模式下，语义与激活按权重 $\alpha \in (0,1)$ 融合；一旦 $s_{\mathrm{meta}}$ 越过阈值 $\tau$，语义分支被\textbf{物理性抑制}，检索键完全由激活特征驱动，强制走向防御 / 越狱相关的黄金安全示例，与表面文本无关。$\tau$ 通过动态 FPR 校准确定（见下文）。
\end{definition}

\noindent 该机制的核心洞察是：激活特征在异常越阈时拥有\textbf{物理开关式的覆盖优先权}，而非与语义分支做梯度式权衡。这保证即使攻击者完美伪装表面文本，只要内部激活异常被检测到，检索就不可能被拉向良性无关示例。图~\ref{fig:jointkey} 展示该路由的完整数据流。

\begin{figure*}[t]
\centering
\fbox{\parbox{0.9\textwidth}{\vspace{16mm}\centering
\textbf{Figure placeholder.}\\
将最终投稿版本的联合键/强制路由示意图放置为 `Figure4.(pdf|png|jpg)` 后，可替换本占位框并恢复 \texttt{\textbackslash includegraphics}。\vspace{16mm}}}
\caption{\textbf{阈值化字典序强制优先级路由（Definition~\ref{mech:override}）。}检索键由语义嵌入 $\mathbf{e}_{\mathrm{sem}}$ 与激活特征 $\mathbf{e}_{\mathrm{act}}$ 构成。门控条件 $s_{\mathrm{meta}} > \tau$ 触发硬覆盖：正常模式下两分支按权重融合；一旦越阈，语义分支被显式抑制（suppressed semantic path），检索键完全由激活特征驱动，并被强制路由至黄金安全示例库。该机制确保即使表面文本被完美伪装，只要内部表示异常被检测到，检索也不能被拉向良性无关示例。}
\label{fig:jointkey}
\end{figure*}

\textbf{校准与检索预算。}$s_{\mathrm{meta}}$ 为连续统计量；固定常数阈值要么过触发（正常用户方差导致误报与检索劫持），要么欠触发（高阶 PBU 潜行）。因此我们\emph{不}手工设定覆盖阈值，而采用\textbf{动态 FPR 校准}：在受控\textbf{良性验证集}（设计者策展、无攻击流量、与测试数据\textbf{严格隔离}）上测量正常运行下 $s_{\mathrm{meta}}$ 的经验分布，选择阈值 $\tau$ 使\textbf{假阳性率}不超过规定水平（如 FPR $= 1\%$）。校准遵循三项操作约束：（i）\textbf{更新周期}——每次基座模型版本更新或部署分布显著偏移时重新校准；（ii）\textbf{数据隔离}——校准集与评测/在线数据严格分离，防止数据泄漏导致阈值虚高；（iii）\textbf{校准失败策略}——若校准集不足或分布异常，默认冻结上一版有效阈值并\emph{保守地提高}覆盖频率（即临时降低 $\tau$），以安全优先。这与 APST 式压力测试哲学一致~\cite{apst}。

\noindent 对于 exemplar 载荷，系统最多向同步路径注入 Top-$K$（默认 $K=2$）个黄金示例；缓存命中时直接复用，缓存未命中时才触发向量库查询。该约束保证同步侧载荷与延迟影响可控，并与 \S\ref{sec:reconstruction} 的部署延迟切分一致。修正仍锚定于防守方定义示例；变体在 Oracle 上保持语义保持~\cite{ControlledMutation_SafetyEval,jade-linguistics}。

\subsection{高密度专家行为迁移与银标准数据生成流水线}
\label{sec:silver}

本节的方法目标不是通用数据增广，而是\textbf{在对抗上下文下把高密度专家安全行为迁移给受控变异器}。黄金数据不能仅是（提示、安全回复）对。实证（如 ADM-ES）表明行为克隆与专家迁移\emph{强烈依赖}信息密度与长度；简短线索效果可忽略~\cite{adm-es-arxiv}。

\textbf{黄金示例最小数据模式（Exemplar Schema）。}我们要求每条黄金示例至少包含以下字段：

\begin{itemize}[nosep,leftmargin=*]
\item \texttt{prompt} / \texttt{context}——触发场景与上下文；
\item \texttt{safe\_response}——专家给出的安全回复；
\item \texttt{rationale}——\textbf{评论级理由}，明确说明该回复为何安全、拒绝/编辑背后的伦理或安全逻辑；
\item \texttt{violation\_category}——对应的违规类别标签（若适用）；
\item \texttt{tool\_permission\_constraints}——工具/权限约束（若涉及工具调用场景）；
\item \texttt{provenance} / \texttt{audit\_source}——数据来源与审计溯源。
\end{itemize}

\noindent 评论级理由（\texttt{rationale}）是与短提示对的核心区别：它使变异器能学习\emph{为什么}某个边界是安全的，而不仅是模仿输出模式。

\textbf{银标准数据生成流水线。}ADM-ES 通过对黄金专家数据的受控变异构建银标数据集~\cite{adm-es-arxiv}；有效性依赖信息密度。流水线检索含评论级理由的示例，定义变异目标，运行 LLM 变异器，在三维质量门控下决定接纳或拒绝变体：

\begin{itemize}[nosep,leftmargin=*]
\item \textbf{语义保持}（Semantic preservation）：变体与原始安全回复的 BERTScore $\ge$ 预设阈值；
\item \textbf{策略一致性}（Policy consistency）：变体不违反原始策略规则（由 DSL 策略引擎或裁判集成自动验证）；
\item \textbf{安全边界保持}（Safety-boundary preservation）：变体不跨越安全边界类别——如原始示例为「拒绝」类，变体不得变为「部分遵从」类。
\end{itemize}

\noindent 前两类由自动审查覆盖；第三类引入\textbf{人工 spot-check}（固定比例抽检），与全文的人类锚定叙事闭合。适应度结合语义相似度与行为一致性；诊断输入推荐图（如 UMAP）以识别失败簇并针对性修复~\cite{adm-es-arxiv}。槽位变异与 CSE~\cite{cse-evocontrol,cse-arxiv} 允许对子能力（如工具路由、拒绝一致性）进行外科式压力测试。

%-------------------------------------------------------------------------------
\section{评估：主张—证据—边界的经验检验}
\label{sec:evaluation}
%-------------------------------------------------------------------------------

本章的目标不是“如何跑实验的说明书”，而是将全文主张组织为\textbf{可快速核验}的证据链：\textbf{主问题（RQ）$\to$ 主结果 $\to$ 机制证据 $\to$ 边界与代价 $\to$ 复现契约}。除非显式标注，所有“可部署结论”仅基于 \textbf{Track A（生产对齐）}；\textbf{Track B（实验室上界）}仅用于组件压力测试与机制验证，\textbf{不得与 Track A 混读或混排}。本章指标与 \S\ref{sec:physics} 的保证—退化边界一一对应。

\subsection{评估问题与证据映射（RQ-first）}
\label{sec:eval-scope}

\noindent 本章按研究问题（RQ）组织结果与证据，而不是按协议组件目录组织。为减少 over-claim 与跨轨道外推，所有 RQ 的结论均显式标注其证据位置与边界入口（\S\ref{sec:physics}）。

\paragraph{RQ1（Track A，可部署联合目标）。}
在生产约束下，统一架构是否在\textbf{安全—效用联合目标}上优于强基线？\;\textbf{证据：}表~\ref{tab:tracka-main-cn}（主表）+\S\ref{sec:eval-tracka-readout}（联合读法）。\;\textbf{边界入口：}\S\ref{sec:physics-claims}--\S\ref{sec:physics-failure}。

\paragraph{RQ2（§\ref{sec:reconstruction} 的 fail-safe 协议）。}
两阶段协议是否经验上满足降级正确性（阻断对抗性合理化）、S1 同步可用性，以及恢复必要性？\;\textbf{证据：}\S\ref{sec:eval-minimal}（E1--E3）。\;\textbf{边界入口：}\S\ref{sec:physics-failure}（可用性反噬/服务补偿）。

\paragraph{RQ3（§\ref{sec:mutation} 的 ORA 闭环）。}
Observe--Route--Adapt 的三要素是否分别贡献可测增益，且 override 与高密度理由在对抗上下文中是必要的？\;\textbf{证据：}\S\ref{sec:eval-minimal}（E4--E6）。\;\textbf{边界入口：}\S\ref{sec:physics-regimes}（观测能力分层）。

\paragraph{RQ4（真实攻击面与分布漂移）。}
在 in-the-wild、PAIR、CRA/HPM 等现实攻击表面下，哪些失败簇主导代价与退化？\;\textbf{证据：}\S\ref{sec:eval-tracka-readout} 与 \S\ref{sec:eval-utility}（漂移、SLA、DoS 切片）。\;\textbf{边界入口：}\S\ref{sec:physics-failure}。

\begin{table}[t]
\centering
\small
\setlength{\tabcolsep}{3pt}
\begin{tabular}{@{}p{0.6cm}p{2.55cm}p{2.85cm}@{}}
\toprule
\textbf{RQ} & \textbf{关键证据（主表/主节）} & \textbf{边界入口（§6）} \\
\midrule
RQ1 & 表~\ref{tab:tracka-main-cn}；\S\ref{sec:eval-tracka-readout} & \S\ref{sec:physics-claims}--\S\ref{sec:physics-failure} \\
RQ2 & E1--E3（\S\ref{sec:eval-minimal}） & \S\ref{sec:physics-failure} \\
RQ3 & E4--E6（\S\ref{sec:eval-minimal}） & \S\ref{sec:physics-regimes} \\
RQ4 & \S\ref{sec:eval-tracka-readout}；\S\ref{sec:eval-utility} & \S\ref{sec:physics-failure} \\
\bottomrule
\end{tabular}
\caption{\textbf{RQ--Evidence Map（压缩版）。}本章所有结论均在表中所示证据范围内解释，禁止跨 Track 外推或混读。}
\label{tab:rq-evidence-map}
\end{table}

\medskip
\noindent\textbf{证据来源与陈述纪律。}本章区分三类陈述：（1）\textbf{形式化定理与解耦安全理论}（不依赖单一实验台）；（2）\textbf{文献中的实证数字}（均标注引用来源，本文不改写为自有复现）；（3）\textbf{本文的评估义务}：对统一架构给出可证伪指标与双轨威胁模型下的完整配方。\textbf{Track~A 参考实测}来自开源 harness（协议 \texttt{paper-eval-5}），汇总于表~\ref{tab:tracka-main-cn}；复现契约见 \S\ref{sec:eval-artifact}（主文压缩版），工程细节下沉至附录（见 \S\ref{sec:appendix-eval-artifact}）。

\subsection{评估双轨：生产对齐与实验室上界}
\label{sec:eval-tracks}

\textbf{Track A（生产对齐）。}攻击者仅有\textbf{黑盒 API}：多轮查询、观测响应与元数据（如拒绝原因类别若暴露）。允许公开集合（如 JailbreakBench 的 JBB-Behaviors~\cite{jailbreakbench}、in-the-wild 提示库~\cite{wild-jailbreak-prompts}、PAIR 生成~\cite{pair-jailbreak}）、多轮抽取脚本、HPM 基准的\textbf{提示级}攻击；\textbf{禁止}假设攻击方可访问 victim 梯度或权重。\textbf{主结论}（可部署性、良性效用、与基线的 ASR 差）仅基于 Track A 报告。

\textbf{Track B（实验室上界）。}在受控环境中允许\textbf{白盒/灰盒}能力：GCG 式与 AutoDAN 式~\cite{gcg,autodan} 梯度或优化攻击、表征对齐 PBU（对隐藏状态施加正则）、surrogate 迁移攻击等，用于\textbf{估计组件与管线在压力下的上界}。Track B 结果\textbf{不得}与 Track A 合并为单一 ASR 表；须在图表中分面板或分行标注。

\subsection{数据集分类：资产类型与交互形态}
\label{sec:eval-taxonomy}

为使基准与解耦安全理论及制品评估（AE）预期对齐，我们将 Track A/B 数据按\textbf{资产类型}（被攻击对象）与\textbf{交互形态}（单轮 vs.\ 有状态；黑盒 vs.\ 白盒）组织。表~\ref{tab:dataset-taxonomy} 列出代表性来源；配套仓库以机器可读 \textbf{YAML} 清单保存相同字段（\S\ref{sec:eval-artifact}）。

\begin{table*}[t]
\centering
\small
\setlength{\tabcolsep}{3pt}
\begin{tabular}{@{}p{2.2cm}p{3.8cm}p{3.0cm}p{4.6cm}@{}}
\toprule
\textbf{资产类型} & \textbf{代表性数据} & \textbf{交互形态} & \textbf{主指标 / 轨道} \\
\midrule
行为约束 & JBB-Behaviors、AdvBench、StrongREJECT~\cite{jailbreakbench,strongreject} & 单轮；misuse + 良性对照 & RSR、ASR、良性 FPR；Track A \\
提示/策略工件 & JailbreakBench artifacts、in-the-wild、PAIR 产物~\cite{jailbreakbench,wild-jailbreak-prompts,pair-jailbreak} & 单轮；分布漂移 & 真实分布下 ASR；Track A \\
隐藏上下文 & CRA 式重构、PBU~\cite{cra-reconstruction,pbu-reconstruct} & 多轮有状态 & 末轮 F1、会话 max-F1；Track A/B \\
系统配置 & SPE-LLM 式抽取~\cite{spe-llm} & 多轮/抽取 & 相对参考提示的 F1；Track A/B \\
心理状态 & HPM / 心理测量套件~\cite{hpm-jailbreak} & 多轮；胁迫 & ASR、压力代理、PCS 类分数；Track A \\
\midrule
\multicolumn{4}{@{}p{0.96\textwidth}@{}}{\textbf{白盒上界（Track B）：}梯度类方法（GCG~\cite{gcg}、AutoDAN~\cite{autodan}）在 AdvBench 兼容目标上探测上界；结果\textbf{不得}与 Track A 主表 ASR 混排。} \\
\bottomrule
\end{tabular}
\caption{\textbf{评估覆盖矩阵（非清单）。}资产类型 $\times$ 交互形态及代表性社区基准，用于定义第 5 章的覆盖面契约：每一行同时指明所属轨道（Track A/B）与主报告指标，从而支撑后文对可部署性（RQ1）、协议证据（RQ2）、ORA 机制证据（RQ3）与漂移/代价边界（RQ4）的经验检验。\texttt{Decoupled-LLM-Gateway} 中 YAML 清单与此表一致，便于 AE。}
\label{tab:dataset-taxonomy}
\end{table*}

\subsection{攻击基线：复现配方与披露等级}
\label{sec:eval-attack}

我们对统一架构进行下列压力测试；每条须固定\textbf{版本、预算与成功判据}以便复现。

\textbf{（1）基于梯度的越狱（仅 Track B）。}须同时报告 \textbf{GCG} 式后缀/前缀优化~\cite{gcg} 与 \textbf{AutoDAN}~\cite{autodan}（语义上接近人类、可绕过困惑度检测与纯表面防御）。对每种方法写明 victim 与 surrogate（若迁移）、对抗串长度、迭代步数与早停规则；成功判据为（i）有害行为分类器触发，或（ii）与参考有害完成串的匹配率，二者取协议预先规定者。

\textbf{（2）多轮智能体抽取（Track A/B）。}采用 \textbf{Conversation Reconstruction Attack（CRA）} 仓库格式~\cite{cra-reconstruction} 或等价结构化历史进行多轮重构与 PBU 式探针；在 Mastermind 等智能体框架中\textbf{选定其一为主基准}并说明取舍（交互轮式、是否假设部分系统提示可见）。固定\textbf{最大轮数} $N_{\max}$、每轮 Token 上限与总预算；报告达到目标保真度（如 F1 $>\tau$）的\textbf{轮数分布}、\textbf{末轮 F1} 与\textbf{会话内 max-F1}，而非仅最终轮。

\textbf{（3）心理学操纵（HPM，Track A）。}采用将操纵刻画为跨轮\textbf{心理状态/策略性脆弱性资产}的 HPM 基准~\cite{hpm-jailbreak}；写明是否扩展语言、判定为\textbf{自动分类器}还是\textbf{人工裁定}。若含压迫性措辞，须在伦理与内容警告层面按机构政策处理。

\textbf{（4）自适应攻击者与披露等级。}将「知晓防御」操作化为\textbf{L0--L2}：L0 仅知存在护栏；L1 知模块名称与接口；L2 知伪代码级流程与关键超参（不含我方密钥与内部随机源）。诱饵 DoS 与表征对齐 PBU 须注明所属披露等级与 Track。

\textbf{（5）诱饵 DoS（Track A）。}攻击目标为耗尽\textbf{时空轮次预算}而非内容抽取；报告同节点\textbf{良性 SLA}：可用请求比例、p95 延迟（\S\ref{sec:eval-utility}）。

\textbf{（6）表征对齐 PBU（仅 Track B）。}攻击者联合优化表面文本与隐藏状态，使轨迹接近良性黄金基线。须写明与监控器\textbf{是否同构模型}、优化器步数及是否允许查询 victim 隐藏状态。主张轮次调度与分散增加双重对齐难度时，在本基线下给出 ASR 与\textbf{攻击成本曲线}。

\textbf{（7）标准有害单轮集（Track A）。}优先采用 JailbreakBench 的 \textbf{JBB-Behaviors}~\cite{jailbreakbench}（\textbf{100 条 misuse} 与 \textbf{100 条 benign}，与策略对齐的类别），使 RSR 与\textbf{良性对照集上的误拒率}联合报告，避免单指标刷分。保留与 GCG/AutoDAN 复现兼容的 \textbf{AdvBench 型}目标~\cite{gcg,autodan}。提示表版本与哈希写入制品元数据与附录。拒绝判定可采用\textbf{启发式}（低成本回归）或\textbf{外置 HTTP 裁判服务}（主文审计对齐推荐）；二者须在附录披露协议与一致性（如子集人工 $\kappa$）。

\textbf{（8）PAIR 与 in-the-wild 提示（Track A）。}纳入\textbf{PAIR}~\cite{pair-jailbreak} 作为低查询成本的黑盒语义攻击，直接压测 \S\ref{sec:paradigm-economics} 中的信息经济学界 $C_{\mathrm{attack}} > V(t)$。纳入大规模 \textbf{in-the-wild} 越狱提示集~\cite{wild-jailbreak-prompts}，以贴近真实世界提示分布、超越纯学术列表。

\subsection{防御基线}
\label{sec:eval-defense}

\textbf{统一架构}（上下文混淆、带分布评估的混合裁判集成、防御性降级、元认知监控器、联合键检索、高密度专家迁移）与下列\textbf{防御基线}在\textbf{同一良性集、同一攻击预算}下对比；所有基线共享相同基座 LLM 与解码设置，除非消融另有说明。

\textbf{（1）仅单 guard：}单一现成安全模型（如 Llama-Guard）为唯一裁判，无上下文混淆、无分布评估、无联合键。

\textbf{（2）强系统提示 + 单 guard：}在（1）之上叠加\textbf{固定强系统提示}与标准聊天模板，代表常见生产「轻量组合」。

\textbf{（3）RAG 安全检索（无联合键）：}向模型注入 Top-$K$ 安全示例，但检索键\textbf{仅为}提示语义嵌入——用于隔离「联合键」贡献。

\textbf{（4）输入扰动：}SmoothLLM 式随机化/扰动，无架构重建或受控变异。

\textbf{（5）仅意图过滤：}语义或意图防火墙，无 ID 对齐或诱饵。

\textbf{（6）商品化 SOTA 安全分类器（冻结外置 guard）。}近年越狱与红队基准常将\textbf{强现成的输入/输出安全模型}作为可比的「SOTA guard」基线。本文将基线（6）\textbf{明确定位于此角色}：以\textbf{固定命名版本}的公开分类器（如 Llama Guard~\cite{llama-guard}；revision 写入制品元数据），经 OpenAI 兼容的 HTTP 裁判或输出守卫路径（\S\ref{sec:eval-artifact}）审核相同提示与完成体，并与其它行保持数据集、预算、解码与基座对话模型一致，从而将\textbf{以网关为中心的架构}（上下文混淆、诱饵、降级路由）与\textbf{独立分类器栈}区分开来。在该端点启用时，结果以\textbf{附录或扩展表}形式给出，与表~\ref{tab:tracka-main-cn} 中九种网关 \texttt{X-Gateway-Experiment-Mode} 基线并列。

\noindent\textbf{与领域现状的关系。}LLM 安全并无单一公认的「唯一 SOTA」：并行路线包括冻结版审核器、针对不可信片段的分隔符/结构化包装策略（文中 \texttt{structured\_wrap} 为轻量代理）、推理时扰动（SmoothLLM 式，对应基线（4））以及训练期对齐等。本文以\textbf{多条可复现基线构成的向量}——（1）--（6）及网关消融——进行对比，而非单一排行榜标量。探索性实验可增列其他开放审核模型或厂商 API，但须在附录写明版本号与裁判提示模板。

\subsection{数据集划分与信息泄漏控制}
\label{sec:eval-splits}

\textbf{不相交集合。}规范评测集、变异器生成的变体集、策略/阈值调参集\textbf{三者互不相交}。用于联合键与元认知阈值选择的\textbf{良性验证集}须与最终\textbf{冻结测试集}独立。

\textbf{划分流程。}\textbf{训练/调参}：攻击轨迹与失败模式仅用于候选策略与阈值搜索；\textbf{验证}：在 hold-out 攻击分布上选阈值（如 FPR 目标）；\textbf{测试}：冻结全部超参后一次性报告 ASR、残余抽取与效用。Oracle 标注若参与变异器训练，须在文中声明并考虑\textbf{泄漏风险}与缓解（如仅使用脱敏 Oracle）。

\subsection{指标的操作定义}
\label{sec:eval-metrics}

\textbf{安全与拒绝。}\textbf{RSR}（Refusal Success Rate）：在\textbf{有害/违规提示集}上，模型输出被裁判或协议判定为\textbf{安全拒绝}（含模板降级输出）的比例，按\textbf{请求级}计数。\textbf{ASR}（Attack Success Rate）：攻击目标达成比例，定义与数据集一致（如越狱成功生成有害内容、或抽取 F1 超阈）。\textbf{有害性与空越狱。}判定应遵循 \textbf{StrongREJECT 式}协议~\cite{strongreject}（或等效）：仅当完成体按基准标准\emph{确实}有害时才计为越狱成功，不得将无实质拒绝但输出废话的\textbf{空越狱}算作成功——以免 ASR 相对人工判断虚高。

\textbf{残余抽取。}多轮设定下报告\textbf{最后一轮}与\textbf{会话内最大 F1}（相对嵌入或人工标注的参考秘密串）；二者均需说明参考串构造方式（合成秘密、分段哈希等）。

\textbf{代价不等式 $C_{\mathrm{attack}} > V(t)$ 的操作化。}实验中采用\textbf{可计算代理}：令 $C_{\mathrm{attack}}$ 为达到 F1 $>\tau$ 的\textbf{期望查询代价}，$\hat{C} = \mathbb{E}[N_{\mathrm{rounds}}] \cdot c_{\mathrm{token}} + c_{\mathrm{wall}}$（$c_{\mathrm{token}}$、$c_{\mathrm{wall}}$ 为预先规定的单位成本）；令 $V(t)$ 为管理策略给定的\textbf{轮换窗口}内允许的最大可承受抽取损失（或等效轮数上界）。主文可报告\textbf{曲线}：F1--轮数、$\hat{C}$--$\tau$，而非单一不等式数值。

\textbf{裁判 FPR。}在\textbf{held-out 良性请求集}（与攻击集独立）上，安全模块误判为违规或触发不当降级的比例；须注明该良性集的来源与规模。

\textbf{分布评估与策略更新。}每变体进行 $n$ 次随机采样（固定 $n\ge 5$），估计经验失败率 $\hat{p}$；仅当 $\hat{p}$ 在单侧置信水平 $\alpha$ 下显著高于更新阈值时触发策略变更（如 Clopper--Pearson 精确区间或 bootstrap）。禁止以单次贪婪解码定性。

\textbf{运行成本与延迟。}\textbf{端到端变异延迟}分解为：检测、检索、变异器、发布/回滚；各报告 p50/p95/p99。\textbf{同步路径延迟}在声明硬件上测量上下文混淆与轻量元认知的 p99，验证 $<$50\,ms 预算是否满足。\textbf{Token 开销比：}相对无检索基线的额外 prompt Token。

\textbf{审计。}发布门控的溯源与证据链（DUC/DUA-E 风格），关联策略版本与评估运行 ID。

\subsection{统计报告、主指标与多重比较}
\label{sec:eval-stats}

\textbf{随机性。}对 ASR、RSR、FPR、F1 等报告至少 \textbf{3--5 个随机种子}或 bootstrap 95\% 置信区间；解码温度、裁判采样与变异器若涉随机，须固定种子或分层报告。

\textbf{主指标。}部署场景建议以\textbf{联合目标}为主：在 Track A 上同时报告 \textbf{ASR（越低越好）}与\textbf{良性任务成功率 / 不当拒绝率（越高/越低越好）}，避免单指标「刷分」。对归档 harness 主表的具体读法见 \S\ref{sec:eval-tracka-readout}。

\textbf{多重比较。}多攻击、多防御联合展示时，标明\textbf{探索性分析}或采用适当的多重比较说明，主表聚焦预先登记的\textbf{主假设对比}。

\begin{table}[t]
\centering
\small
\setlength{\tabcolsep}{3pt}
\begin{tabular}{@{}p{1.55cm}p{2.05cm}p{2.65cm}@{}}
\toprule
\textbf{指标} & \textbf{用途} & \textbf{主文位置} \\
\midrule
ASR / RSR & 安全主指标（有害集） & 表~\ref{tab:tracka-main-cn}、\S\ref{sec:eval-tracka-readout} \\
良性任务成功率 / 不当拒绝率 & 效用与误拒代价 & 表~\ref{tab:tracka-main-cn}、\S\ref{sec:eval-utility} \\
残余抽取 $F_1$、max-$F_1$ & 抽取边界与退化 & \S\ref{sec:eval-metrics}、\S\ref{sec:eval-minimal} \\
$\hat{C}$--$\tau$ / 轮数曲线 & 代价代理与经济学约束 & \S\ref{sec:eval-metrics}、\S\ref{sec:eval-attack} \\
p95/p99 延迟、可用率 & SLA / S1 可用性 & \S\ref{sec:eval-utility}、E2（\S\ref{sec:eval-minimal}） \\
\bottomrule
\end{tabular}
\caption{\textbf{Metrics Contract（压缩版）。}主文以联合目标而非单一标量排序；实现细节与脚本路径下沉至附录。}
\label{tab:metrics-contract}
\end{table}

\subsection{Track A 主表与主结果解读（承重墙）}
\label{sec:eval-tracka-main}

\noindent 本节集中回答 RQ1（\S\ref{sec:eval-scope}）：在可部署约束下，统一架构是否在\textbf{安全—效用联合目标}上优于强基线。所有可部署结论仅基于 Track A；Track B 结果\textbf{不得}与本节混读或混排。

\begin{table}[t]
\centering\small
\setlength{\tabcolsep}{4pt}
\begin{tabularx}{\columnwidth}{@{}lccc@{}}
\toprule
\textbf{GW} & \textbf{RSR$_{\mathrm{harm}}$} & \textbf{FPR$_{\mathrm{benign}}$} & \textbf{RSR$_{\mathrm{itw}}$} \\
\midrule
direct\_upstream & 1.0000 & 0.0000 & 1.0000 \\
intent\_only & 1.0000 & 0.0000 & 1.0000 \\
intent\_only\_plus & 1.0000 & 0.0000 & 1.0000 \\
no\_decoy & 1.0000 & 0.0000 & 1.0000 \\
single\_guard\_only & 1.0000 & 0.0000 & 1.0000 \\
strong\_system\_guard & 1.0000 & 0.0000 & 1.0000 \\
structured\_wrap & 1.0000 & 0.0000 & 1.0000 \\
structured\_wrap\_plus & 1.0000 & 0.0000 & 1.0000 \\
unified & 1.0000 & 0.0000 & 1.0000 \\
\bottomrule
\end{tabularx}
\caption{\textbf{Track A 主表（USENIX 友好压缩版，本次复现切片）。}为适配双栏版式，仅展示主读法所需三列：有害集拒绝成功率 RSR、良性误拒率 FPR、in-the-wild RSR。该切片配置：dcbf=v2-proxy-ensemble，judge=heuristic，OG=on。完整逐数据集/逐指标切片见制品导出表。}
\label{tab:tracka-main-cn}
\end{table}



\subsection{如何阅读 Track~A 主表：联合指标与分布漂移}
\label{sec:eval-tracka-readout}

表~\ref{tab:tracka-main-cn} 须按 \S\ref{sec:eval-stats} 的\textbf{联合目标}阅读：有害集上的 \textbf{RSR}（越高表示对违规提示的安全拒绝越多）、\textbf{良性不当拒答率}，以及（若 JSON 导出或附录补充列中包含）\textbf{in-the-wild}、StrongREJECT 式列表上的拒绝率。下列模式在已归档的 \texttt{paper-eval-5} 跑分（含用于诊断的单种子切片\footnote{\textbf{探索性单种子示例}（仅种子 42，\texttt{deepseek-chat}，启发式裁判，制品中一次完整矩阵的 checkpoint JSON）。\textbf{并非}表~\ref{tab:tracka-main-cn} 的跨种子汇总格点；换种子或 HTTP 裁判可显著改变数值。示例比例：\texttt{direct\_upstream} 与 \texttt{unified} 的有害清单 RSR 分别为 $0.85$ 与 $0.91$，良性不当拒答率 $0.05$ 与 $0.16$；\texttt{unified} 在 in-the-wild 上 RSR $0.225$、有害清单 RSR $0.91$；\texttt{structured\_wrap} 有害清单 RSR $0.99$、良性 FPR $0.47$；\texttt{strong\_system\_guard} in-the-wild RSR $0.59$、良性 FPR $0.32$。该切片上，默认 harness 探针下单轮与多轮抽取 $F_1$ 对各防御均为 $0$。}）中反复出现，用于把数字与「解耦安全」叙事对齐---它们\textbf{不是}单一标量排名。

\paragraph{跨路由的安全--效用张力。}
在同一基座模型上，相对 \texttt{direct\_upstream}，启用完整网关（\texttt{unified}）往往同时\textbf{抬高}有害集 RSR 与\textbf{抬高}良性不当拒答。这是部署侧可预期的权衡，而非 harness 故障。网关消融（\texttt{intent\_only}、\texttt{no\_decoy} 等）也不保证在每一列上 Pareto 改进：部分子栈可能降低良性 FPR 同时改变有害 RSR，故须以\textbf{多列向量}比较防御，忌用单一数字排序。

\paragraph{清单有害 vs.\ in-the-wild 漂移。}
制品在\textbf{版本化标准有害行为列表}（JBB-Behaviors 式 misuse 文件）与\textbf{in-the-wild} 越狱提示集上均报告 RSR。两列差距大表明\textbf{分布漂移}：在 curated misuse 上很强的防御，在真实越狱表面形态下拒绝率可能仍显著偏低。关于「可部署鲁棒性」的主张应\textbf{并列引用}两列；仅报 JBB 式 RSR 易\textbf{高估}野外就绪度。

\paragraph{强系统提示与分隔符类基线。}
在启发式裁判下，\texttt{structured\_wrap}、\texttt{strong\_system\_guard} 往往可在有害集上达到很高 RSR，但常伴随明显升高的良性误拒。本文将其理解为\textbf{提示层/格式层硬化并非零成本}：它可能在阻断显式违规的同时，对良性流量过度触发拒绝。褒扬「强系统提示」或分隔符策略时，须以良性列与 HPM 代理列为条件。

\paragraph{代理边界（多轮、HPM、诱饵压测、延迟）。}
默认 harness 模板上单轮与多轮抽取 $F_1$ 归零，是\textbf{该探针}在登记阈值与预算下的结果，\textbf{不}等同于对所有 CRA/SPE 类攻击的免疫（\S\ref{sec:eval-attack}(2)）。HPM 行为\textbf{小样本压力代理}（默认制品文件中为五条提示）；量级宜作示意。诱饵--SLA 场景报告的是并发压测后的\textbf{端到端}良性延迟分位数，\textbf{不是}同步网关微秒级预算是否满足；\S\ref{sec:reconstruction} 中 $<$50\,ms 路径主张须另设仪器测量（\S\ref{sec:eval-metrics}）。

\paragraph{种子与汇总。}
正文主结论应以制品声明的种子集及表~\ref{tab:tracka-main-cn} 的汇总统计为准（见 \S\ref{sec:eval-scope}）；单种子 JSON 用于诊断，\textbf{不可替代}汇总表作论文级断言。

\begin{table}[t]
\centering\small
\setlength{\tabcolsep}{3pt}
\begin{tabularx}{\columnwidth}{@{}lccc@{}}
\toprule
\textbf{GW} & \textbf{FPR$_{\mathrm{benign}}$ (on/off/$\Delta$)} & \textbf{RSR$_{\mathrm{harm}}$ (on/off/$\Delta$)} & \textbf{mean-$F_1$$_{\mathrm{extract}}$ (on/off/$\Delta$)} \\
\midrule
structured\_wrap & 0.0000/0.0000/0.0000 & 1.0000/1.0000/0.0000 & 0.4346/0.4346/0.0000 \\
unified & 0.0000/0.0000/0.0000 & 1.0000/1.0000/0.0000 & 0.4346/0.4346/0.0000 \\
\bottomrule
\end{tabularx}
\caption{\textbf{Table 6：输出守卫（OG）消融（USENIX 友好压缩版，本次复现切片）。}仅展示最关键的效用/安全/抽取三项指标；配置：dcbf=v2-proxy-ensemble，judge=heuristic。}
\label{tab:tracka-guard-cn}
\end{table}



\subsection{边界、负结果与代价汇总（与第 6 章闭环）}
\label{sec:eval-boundary-summary}

\noindent 本节把本章分散的“坏消息”集中呈现，降低“在藏失败”的 attack surface，并将经验退化与 \S\ref{sec:physics-failure} 的退化分层对齐。

\begin{table}[t]
\centering
\small
\setlength{\tabcolsep}{3pt}
\begin{tabular}{@{}p{1.65cm}p{2.0cm}p{1.45cm}p{1.0cm}@{}}
\toprule
\textbf{失败簇} & \textbf{触发条件（示例）} & \textbf{主代价} & \textbf{§6} \\
\midrule
漂移 gap & curated $\to$ in-the-wild 形态变化 & 误拒/漏拒上升 & \S\ref{sec:physics-failure} \\
降级过触发 & 阈值偏保守/压力信号频繁越阈 & 良性可用性下降 & \S\ref{sec:physics-failure} \\
观测不足 & 黑盒/信号稀疏导致监控退化 & 机制增益变弱 & \S\ref{sec:physics-regimes} \\
\bottomrule
\end{tabular}
\caption{\textbf{Boundary Evidence Summary（压缩版，占位）。}列出最主导的失败簇、触发条件与代价，并指向第 6 章对应的退化类别；实证填充应来自主表切片与 E1--E6 的对照结果。}
\label{tab:boundary-evidence-summary}
\end{table}

\paragraph{漂移（drift）。}
curated misuse 与 in-the-wild 的差距是主结果的一部分：若二者差距显著，则任何“可部署鲁棒性”的表述必须带条件，并在 \S\ref{sec:physics-failure} 中作为边界声明，而非在主文中隐含外推。

\paragraph{失败簇（failure clusters）。}
对多轮抽取、CRA/HPM、以及 DoS 压力测试，应报告最主导的失败簇与触发条件（例如：预算不足、裁判不确定区间扩大、或降级触发频率过高），并明确对应的最低退化形式（\S\ref{sec:physics-failure}）。

\paragraph{代价（cost）。}
任何安全增益必须同时报告良性误拒代价与同步 SLA（p95/p99、可用率）；否则会被视为单指标刷分或不可部署。E2（\S\ref{sec:eval-minimal}）是对 S1 的最小经验支撑。

\subsection{可用性、诱饵 DoS 与切片}
\label{sec:eval-utility}

\textbf{良性效用。}除 FPR 外，报告\textbf{代表性良性任务集}上的完成率（如摘要、编码、多轮工具调用），以捕捉「过度拒绝」对业务的损害。

\textbf{诱饵 DoS 与降级 SLA 量化。}在固定速率限制下，模拟高频诱饵命中，报告\textbf{良性请求可用率}与\textbf{级联超时率}（若适用），与 \S\ref{sec:physics} 一致。降级路径（D1--D4）下同步 API 的 p95/p99 延迟与成功返回率须联合报告，以验证 Service Invariant-S1（\S\ref{sec:degrade-actions}）；对应最小实验 E2（\S\ref{sec:eval-minimal}）。

\textbf{切片（可选）。}若声称跨场景部署，按\textbf{语言、领域、请求长度}报告安全--效用权衡的切片表，避免仅报告全局平均。

\subsection{人因与强制人工校准}
\label{sec:eval-human}

\textbf{强制人工校准（如 5\% 抽检）。}协议须包括：标注指南、审核员间一致性（如 Cohen's $\kappa$）、是否与作者\textbf{独立}。校准用于纠正集成 FPR 与零日共识失败，不作为主指标的唯一来源。

\subsection{最小可验证实验包}
\label{sec:eval-minimal}

在篇幅受限时，优先完成支撑三条理论轴的\textbf{最小集合}：（i）\textbf{信息经济学}：多轮抽取 F1--轮数曲线 + 上下文混淆（ID/诱饵）消融；（ii）\textbf{价值锚定}：HPM 集上 CoT 开/关与防御性降级路径的 ASR 与良性成功率；（iii）\textbf{表征异常}：PBU 场景下\textbf{联合键 vs 纯语义检索}的检索命中率与下游修正成功率；表征对齐 PBU 置于 Track B 作为针对性压力测试。

\paragraph{机制证据墙（把实验清单转成可核验结论）。}
为使 §\ref{sec:axiological} 与 §\ref{sec:mutation} 的机制主张在主文中可被快速核验，我们将最小证据包按“结论句—证据句—边界句”写成两组机制证据墙，并在后文给出对应的 E1--E6 细化协议。

\paragraph{Fail-safe（对应 §\ref{sec:axiological}；E1--E3）。}
\textbf{结论句（D/R）。}在触发 T1/T2/T3 的对抗场景下，两阶段 fail-safe 的降级路径在经验上阻断对抗性合理化并维持同步可用性（S1），且恢复通道对合法边缘案例是必要的。%
\textbf{证据句。}该结论由 E1（降级正确性）、E2（S1 的 p95/p99 与成功返回率）、E3（w/ vs.\ w/o 恢复的边缘案例处理率）在同一评估合同下支撑。%
\textbf{边界句。}当部署不满足观测/审计或恢复层级能力（\S\ref{sec:physics-assumptions}）时，上述结论退化为“可拒绝但难恢复”的服务补偿（\S\ref{sec:physics-failure}）。

\paragraph{ORA（对应 §\ref{sec:mutation}；E4--E6）。}
\textbf{结论句（D/R）。}在表面良性但内部异常的对抗会话中，Observe--Route--Adapt 闭环的监控、override 与高密度理由三要素分别贡献可测增益，且 override 在越阈时优于简单融合。%
\textbf{证据句。}该结论由 E4（监控 TPR/FPR 与 reset 后 unsafe continuation 降幅）、E5（三种检索策略对检索命中率与下游修正成功率的对照）、E6（三种理由密度对银标质量门控的对照）支撑。%
\textbf{边界句。}在完全黑盒、无法获得足够表示级信号的 regime 下（\S\ref{sec:physics-regimes}），监控与路由只能退化为输出级代理，其增益与代价需单独报告，不得与白盒/半白盒结论混称同一保证。

\textbf{两阶段 Fail-Safe 协议专属最小证据包。}为使 \S\ref{sec:axiological} 的协议定义可证伪，我们在上述三轴之外补充以下三组实验（E1--E3），每组直接对应协议术语：

\begin{itemize}[nosep,leftmargin=*]
\item \textbf{E1：降级正确性（对应 T1/T2/T3 $\to$ D1--D4）。}在 HPM / 公理冲突类攻击场景下，比较启用 D1--D4 降级路径与仅使用强拒绝模板（无能力隔离）时的 \emph{unsafe continuation rate}（模型在降级后仍产出不安全内容的比例）与 \emph{adversarial rationalization rate}（模型为顺从恶意指令编造合理化论据的比例）。E1 回答：降级是否真的阻断了对抗性合理化？RSR 与良性成功率联合报告。
\item \textbf{E2：可用性与 S1（对应 Service Invariant-S1）。}在触发降级的压力场景下，测量同步路径的\textbf{成功返回率}、\textbf{p95/p99 端到端延迟}与\textbf{级联超时/业务失败率}。对比基线为：（a）无降级直通上游；（b）简单挂起等待人工审核。E2 直接支撑 S1——证明降级路径在有界延迟内保持同步 API 可用。
\item \textbf{E3：恢复必要性（对应 Tier-R1/R2）。}对比仅 D1 refusal（无恢复）与 D1/D2 + 异构恢复（至少含 Tier-R1 规则化可判定恢复）时的\textbf{合法边缘案例处理率}（edge-case recovery rate）——即因安全模糊性被错误降级的合法请求中，经恢复后成功服务的比例——与整体安全指标（误放行率）。E3 回答：不做恢复会损失多少可接受任务？加入恢复是否在不放松安全的前提下改善了边缘案例处理？
\end{itemize}

\noindent E1--E3 的指标操作定义与 \S\ref{sec:eval-metrics} 一致；数据来源为 Track A HPM 代理集、良性任务集与合法边缘案例子集（配套仓库提供提示文件与评分脚本）。

\textbf{受控变异专属最小证据包。}为使 \S\ref{sec:mutation} 的 Observe--Route--Adapt 闭环可证伪，我们补充以下三组实验（E4--E6），每组直接对应第 4 章的方法定义：

\begin{itemize}[nosep,leftmargin=*]
\item \textbf{E4：元认知监控有效性（对应 Observe / $s_{\mathrm{meta}}$）。}在``表面良性、内部异常''的对抗会话（如完美 PBU、任务重构伪装）上，报告异常分数 $s_{\mathrm{meta}}$ 的\textbf{良性 FPR}、\textbf{恶意 TPR}；在 reset 触发后，测量 \emph{unsafe continuation rate} 的下降幅度。E4 回答：元认知监控器是否真的能把表面良性但内部异常的会话拉出来？
\item \textbf{E5：Override 机制必要性（对应 Route / Definition~\ref{mech:override}）。}在 PBU / task-reframing 场景下，直接对比三种检索策略：（a）semantic-only retrieval；（b）weighted fusion（$\alpha$-加权拼接）；（c）thresholded lexicographic override。指标为\textbf{检索命中率}（retrieval hit rate：Top-$K$ 中安全相关示例的比例）与\textbf{下游修正成功率}（downstream correction success：注入示例后变异器产出安全变体的比例）。E5 回答：为什么不能用简单融合，必须用字典序覆盖？
\item \textbf{E6：高密度理由价值（对应 Adapt / \S\ref{sec:silver}）。}对比三种黄金示例密度级别：（a）prompt-response pairs only；（b）prompt-response + short rationale；（c）prompt-response + comment-level rationale。指标与 \S\ref{sec:silver} 的三维质量门控对齐：\textbf{语义保持}（BERTScore $\ge$ 阈值）、\textbf{策略一致性}与\textbf{安全边界保持}。E6 回答：评论级理由是否真的比短线索 exemplar 有效？
\end{itemize}

\noindent E4--E6 的数据来源为 Track A PBU/HPM 代理集与银标准流水线输出；指标操作定义与 \S\ref{sec:eval-metrics} 一致。

\subsection{参考实现与实证验证}
\label{sec:eval-artifact}

我们发布参考实现与评估 harness，使 Track A 的主表结果可由固定提示文件版本/哈希、随机种子与表格导出脚本复现。Track A 的汇总主表由脚本从归档 JSON 重建；提示文件哈希与评估矩阵见附录 \S\ref{sec:appendix-eval}。Track B 的白盒上界与部分压力测试不作为制品承诺范围（与 \S\ref{sec:eval-tracks} 一致）。

为降低主文 attack surface，本节不在正文展开端点、环境变量、脚本与目录等工程细节；完整复现与制品细节见附录 \S\ref{sec:appendix-eval-artifact}。

\subsection{反馈环与治理}

重建消费攻击轨迹与失败模式；变异以受控方式产生它们。示例引导变异器生成压迫当前策略的对抗变体；重建层响应并更新；元认知监控器在检测到不稳定时可触发对抗性重置。轮次调度与变体分散保护基准资产~\cite{ControlledMutation_SafetyEval}。强制人工校准量化并纠正集成 FPR~\cite{tc260-governance}。证据链将每次评估与变异器步骤、调度器状态及策略版本关联，用于发布门控与 attestation~\cite{dual-use-attestations,ControlledMutation_SafetyEval}。

\subsection{小结：评估如何支撑全文 claim}
\label{sec:eval-summary}

本章以 Track A 的主表（表~\ref{tab:tracka-main-cn}）回答“可部署联合目标下是否优于强基线”（RQ1），以 E1--E6（\S\ref{sec:eval-minimal}）验证 fail-safe 协议与 ORA 闭环的机制必要性（RQ2--RQ3），并通过漂移、SLA 与 DoS 切片（\S\ref{sec:eval-tracka-readout}、\S\ref{sec:eval-utility}）将经验代价与 \S\ref{sec:physics} 的退化边界闭环对齐（RQ4）。除非显式标注为 Track B 上界，本章不对评估合同外攻击面做外推。

%-------------------------------------------------------------------------------
\section{解耦安全的边界条件与失效物理学}
\label{sec:physics}
%-------------------------------------------------------------------------------

本章不再引入新的机制，而是对全文主张进行边界化收束。我们的目标不是再次解释理论、系统或评估细节，而是回答四个问题：本文到底主张什么，不主张什么；这些主张依赖哪些前提；当前提失效时系统如何退化；以及为何这些有限主张仍然具有部署意义。因而，本章应被理解为全文的 \emph{boundary of validity}，而非系统实现或复现实务的补充说明。

\subsection{不保证与所主张的保证}
\label{sec:physics-claims}

\noindent\textit{本节回答：我们到底 claim 什么，不 claim 什么？}

\textbf{Non-goals（非目标）。}本文不主张全局意义上的通用安全证明，也不主张存在对自然语言开放交互可一劳永逸适用的完备安全判定。本文同样不将 RBAC、信息流控制（IFC）或组织治理视为可由运行时机制替代的外围条件；对于静态或永久高价值资产，它们构成独立且必要的硬边界。最后，本文不声称可在无人类锚定的前提下实现完全机器自治的安全闭环。

\textbf{Guarantees（有界主张）。}本文主张的是一组有界、外部化、可审计的性质，而非全局安全结论。在 \S\ref{sec:physics-assumptions} 所述前提成立时，上述性质具体化为：其一，安全义务可以从不可判定的全局隐式空间退回到实例级运行时断言；其二，当高风险操纵态出现时，系统可以通过两阶段 fail-safe 协议收敛到受限执行模式，并保留审计链与恢复路径；其三，部署结论必须按白盒、半白盒与黑盒能力分层理解，不能跨部署条件外推；其四，带外 Observe--Route--Adapt 闭环只向同步路径输出受限控制信号，而不承担无限制的自主安全裁决。

\noindent 上述主张以 \S\ref{sec:physics-assumptions} 所列前提为条件；前提失效时，其有效范围按 \S\ref{sec:physics-failure} 退化。

\subsection{前提假设}
\label{sec:physics-assumptions}

\noindent\textit{本节回答：上述 claim 靠什么前提成立？}

\noindent 本文的所有边界性主张同时依赖理论、部署与治理三类前提。

\noindent\textbf{理论前提（theoretical preconditions）。}本文关于 DCBF 类前向不变性与轨迹级安全监控的主张，条件于一个核心前提：存在可计算的代理映射 $\phi$，使投影后的低维轨迹对部署者所声明的安全违规具有可辩护的语义对应。该前提并不意味着系统直接验证了自然语言语义本身；相反，它意味着相关理论结论只在代理足够忠实时成立。若该对应关系失真，则这些结论应被理解为 proxy-level monitoring heuristic，而不是对自然语言安全的无条件定理。

\smallskip
\noindent\textbf{部署前提（deployment preconditions）。}本文的运行时主张还依赖两个部署条件。第一，系统必须拥有足以支撑监控与断言的可观测通道；观测能力越弱，系统可维持的安全粒度越低。第二，系统必须能够在同步路径上执行受限降级，并在带外保留恢复与人工校准机制。前者决定了系统能否实施细粒度断言，后者决定了 fail-safe 是否真正构成一个闭合、可追责的服务协议。若二者任一缺失，运行时主张即退化为更弱的输出级过滤或更高频的人类补偿。

\smallskip
\noindent\textbf{治理前提（governance preconditions）。}本文关于“攻击成本超过资产生命周期价值”的经济学主张，条件于资产具有有界、可轮换的生命周期。对静态或永久高价值资产，例如核心 IP、长期敏感数据或不可替代的对齐语料，该前提并不成立；此时，RBAC、信息流控制（IFC）、审计与元数据治理不是可选补充，而是独立且必要的硬边界。换言之，运行时安全机制只能在这些治理边界之内承担动态防御职责，而不能替代资产治理本身。

\smallskip
\noindent 表~\ref{tab:assumption-matrix} 将三类前提压缩为一页式判定视图。该表\textbf{不是}对定理层级的重新证明，而是将本章 contract 条件压缩为可单独阅读的判定视图：回答``在什么前提下可主张什么、前提失效时最低退化到哪里''，而非``理论何以在一般意义下成立''。

\begin{table*}[t]
\centering\small
\setlength{\tabcolsep}{4pt}
\caption[Claim discipline 矩阵]{Claim discipline 矩阵——前提成立时可主张的性质与失效时的最低退化形式。\textit{注：}本表不引入新的理论或实验结果；左列定义前提类别，中列为前提成立时论文可主张的性质，右列为前提失效时系统退化到的最低保证形式。}
\label{tab:assumption-matrix}
\begin{tabularx}{\textwidth}{@{}>{\raggedright\arraybackslash}p{2.35cm} X X@{}}
\toprule
\textbf{前提类别} & \textbf{成立时可主张的性质} & \textbf{失效时的最低退化形式} \\
\midrule
理论级：表征前提（$\phi$ 忠实） & 可主张条件化不变性与轨迹级监控 & 退化为代理启发式 \\
\addlinespace
治理级：资产轮换 + RBAC/IFC + 监督锚定 & 可主张经济学防御与硬边界 & 退化为治理补偿 \\
\addlinespace
部署级：观测前提 & 可主张细粒度干预（依 regime） & 退化为块级后置断言 \\
\addlinespace
部署级：S1 与异构恢复 & 可主张完整降级—恢复链 & 退化为服务补偿 \\
\bottomrule
\end{tabularx}
\end{table*}

\subsection{失效模式与退化}
\label{sec:physics-failure}

\noindent\textit{本节回答：前提失效时系统如何退化？}

\noindent 当前提失效时，系统并非二值崩塌，而是按不同路径退化到不同的最低保证级别。

本文的失效物理学不是``哪里会坏''的清单，而是``当前提失效时，系统会退化到何种最低保证''的分层描述。下文将退化归纳为四类，每一类都对应一个不同的剩余安全级别。

\textbf{类别 1：经济学前提失效。}当经济学前提失效（资产不可轮换、泄漏后不可回收或其价值远高于可实施查询成本）时，信息经济学上的防御充分性类主张不再成立；最低剩余保证是运行时断言与审计链仍可保留，且安全义务须转移至 RBAC、IFC 与数据治理硬边界。

\textbf{类别 2：监控前提失效。}当监控前提失效（代理映射 $\phi$ 或元认知监控器不忠实于部署者关心的违规模式）时，DCBF/轨迹级强定理类主张不再成立；最低剩余保证是 fail-safe 降级与审计追溯仍在，但主动检测与轨迹级区分显著削弱（典型来源包括嵌入不敏感、基线失真或代理空间不足）。

\textbf{类别 3：可用性反噬。}当可用性前提失效（外部强制约束引发反噬，含算法级拒绝服务与边界信号泄露等张力）时，同时主张高安全与高可用不再成立；最低剩余保证是安全底线仍可维持，但服务连续性、边缘可用性与误拒绝率会恶化。相关张力须由 \S\ref{sec:eval-utility} 的 SLA、降级与恢复实验显式量化。

\textbf{类别 4：未覆盖攻击面。}当攻击面超出本文运行时覆盖（如隐蔽信道、多语言混淆、因果有害输出、训练后门、社会技术长尾等）时，闭合性主张不再成立；最低剩余保证是外围审计、流程治理与外部检测的配合作用（并入主张属范畴性 over-claim，故 out-of-scope）。

\smallskip
\noindent 本文并不将上述退化视为附带限制，而将其视为主张成立条件的一部分。换言之，我们的目标不是否认这些失效路径的存在，而是显式说明：当前提失效时，哪些主张不再成立，以及系统退化到何种最低保证级别。正因如此，本文不把运行时安全框架表述为可脱离治理、访问控制或人类锚定而独立完备的终局方案，而将其定位为一组在明确部署条件下可被外部核验、可被审计追溯、并在失败时可收敛的有界契约。若这些边界不被显式陈述，则本文更容易被误读为对自然语言安全、内生对齐或广义鲁棒性的强承诺，而这并非当前版本能够由理论与评估共同支持的结论。

\subsection{按部署能力分层的保证}
\label{sec:physics-regimes}

\noindent\textit{本节回答：不同部署能力下保证边界如何分层？}

\noindent 不同部署能力仅支持不同粒度的安全主张，因此 guarantee 必须按 regime 分层理解。

\textbf{白盒部署。}当系统可获得细粒度生成信号并能在推理路径上施加约束时，实例级断言可以锚定于更细粒度的生成状态；在这一条件下，本文关于运行时投影与强干预粒度的主张最为充分。即便如此，这些主张仍然条件于 \S\ref{sec:physics-assumptions} 所述代理忠实性与服务前提，而不应被读作无条件的全局安全结论。

\textbf{半白盒部署。}当系统只能获得部分内部统计信号、但无法直接控制生成路径时，其可维持的保证介于细粒度运行时约束与输出级后置断言之间。此时，系统仍可利用异常统计实施监控与受限降级，但不能主张原位的 token 级修正；因此，这一层的结论应被理解为``可检测、可收敛''，而非``可逐步纠偏''。

\textbf{黑盒闭源部署。}当系统只能访问输入、完整输出及极少量外显元信号时，本文不再主张细粒度运行时干预；系统至多维持块级后置断言、外部 fail-safe 路由与审计追溯。此时，系统仍可提供 containment 与治理接口，但不能将黑盒场景下的输出级控制与白盒场景下的细粒度安全主张混称为同一保证。

\textbf{表征监控边界。}本文从未主张复杂社会学构念在表示空间中天然可分。相反，\S\ref{sec:metacog} 采用的是相对良性基线的异常检测代理，其价值在于提供一个可部署、可审计的偏离信号，而不是字面意义上的``认知状态识别器''。因此，本章关于表示级监控的主张应被理解为异常检测意义上的 \emph{bounded evidence}，而非对内部心理状态的本体论判断。

\subsection{有限保证的部署意义}

\noindent\textit{本节回答：为什么这些有限保证仍有部署意义？}

\noindent 这些有限保证的价值，不在于证明模型内生安全，而在于将安全义务转写为可治理的外部契约。

这些限制并不削弱本文的部署价值；相反，它们界定了本文真正解决的问题：不是证明模型本身具备稳定且普适的内生安全性，而是将安全义务转写为外部、显式、可审计、可收敛的运行时契约。其价值不在于``绝对防止所有攻击''，而在于在明确威胁模型与预算下压缩事故半径、抬高攻击成本，并把不可消除的判断残差转移到可治理的人类与制度接口之上。

本文框架不是脱离 RBAC、IFC 与组织治理的独立终局方案，而与三者互补分层：RBAC/IFC 提供高价值资产访问边界；运行时协议与 ORA 闭环承担动态对抗中的即时安全义务；组织治理与人工审计负责残余不确定性的最终锚定。任何单层的``完备性''主张都构成范畴性越界。

\subsection{小结}

综上，本文中的解耦安全应被理解为一种 \emph{bounded, auditable, deployment-oriented safety contract}：它不承诺通用智能的全局安全证明，而只在显式前提与部署约束下主张实例级运行时断言、fail-safe containment、审计可追溯性与分层退化。其意义不在于消除人类判断，而在于把人类判断从高频、实时、不可追责的关键路径中移出，并重置到低频、可延迟、可治理的制度接口之中。因而，本文应被理解为一种面向部署的边界化安全契约，而非对通用安全智能的全局证明。

%-------------------------------------------------------------------------------
\section{相关工作}
%-------------------------------------------------------------------------------

LLM 数据抽取与重建已有大量研究~\cite{llm-reconstruction-emergent,split-ft-attack,active-reconstruction}；我们通过结构变异与调度而非仅噪声缓解。上下文构建防御在基于意图的防火墙基础上扩展以抵御任务重构~\cite{agentos-arxiv}；元认知监控与示例引导变异强化对 HPM 的防御；混合裁判集成与强制人工校准应对单一裁判偏差。\textbf{基准与社区工件。}越狱评测已围绕 \textbf{JailbreakBench}（JBB-Behaviors；misuse/良性对照）~\cite{jailbreakbench}、抑制空越狱的 \textbf{StrongREJECT} 评估器~\cite{strongreject}、低查询黑盒攻击 \textbf{PAIR}~\cite{pair-jailbreak} 与大规模 \textbf{in-the-wild} 提示集~\cite{wild-jailbreak-prompts} 收敛；多轮重构遵循 CRA 类协议~\cite{cra-reconstruction}，系统提示抽取在 SPE-LLM~\cite{spe-llm} 中标准化。Track B 白盒攻击含 GCG~\cite{gcg} 与 AutoDAN~\cite{autodan}。JADE 及语言学平台~\cite{jade-linguistics,jade-db} 为变异策略提供参考；TC260 与 attestation 框架支撑治理~\cite{tc260-governance,tc260-training-data,dual-use-attestations}。SISF~\cite{sisf-arxiv}、AgentOS~\cite{agentos-arxiv}、ADM-ES~\cite{adm-es-arxiv}、CSE~\cite{cse-arxiv} 及价值锚定工作~\cite{epistemic-grounding,hpm-jailbreak} 是我们置于解耦安全理论下统一的构件；ControlledMutation\_SafetyEval~\cite{ControlledMutation_SafetyEval} 提供变异器与数据工厂设计。

\textbf{本文与先前工作的 Delta 贡献。}上述引用中，SISF、AgentOS、ADM-ES、CSE 与 ControlledMutation\_SafetyEval 来自同一研究方向。需明确区分：\textbf{(a)}~SISF 提出了自改进安全框架的闭环概念，但未给出可计算性理论基础或形式化安全代数——本文的 Theorem~1--3 与 Risk poset 框架是全新贡献。\textbf{(b)}~ControlledMutation\_SafetyEval 设计了变异器与数据工厂，但未涉及运行时安全内核或 DCBF 行为监控——本文的 Ring-0 解耦架构、Axiom Hive 投影与 Judge Ensemble 置信度加权是全新系统设计。\textbf{(c)}~AgentOS 提出了意图防火墙，本文的上下文混淆引擎（ID 对齐、诱饵注入）是在此基础上的显著扩展。\textbf{(d)}~本文的核心范式贡献——将 LLM 安全从``隐式对齐''重新定位为``解耦安全的计算必然性''——在上述任何先前工作中均未出现。

%-------------------------------------------------------------------------------
\section{讨论}
%-------------------------------------------------------------------------------

\textbf{Track A 与 Track B 的分轨解读指导。}本文评估协议严格区分\textbf{Track A（生产对齐轨道）}与\textbf{Track B（实验室上界轨道）}。Track A 限定为黑盒查询、公开越狱提示与多轮攻击框架，其结果直接对应「可上线部署」的防御效果；报告指标包括 ASR（攻击成功率）、RSR（拒绝成功率）、良性任务成功率与良性 SLA（p95 延迟）。Track B 额外允许白盒梯度攻击（GCG、AutoDAN）与表征对齐 PBU，其结果仅用于\textbf{探测组件压力上界}——即在最强攻击者假设下的防御天花板。\textbf{关键阅读原则}：Track B 的 ASR 数值\emph{不应}与 Track A 混列或直接比较；二者的攻击者能力假设不同，混排会误导部署决策。Table~\ref{tab:eval-matrix} 完整列出了两轨的攻击面、数据预算与主指标对照。

\textbf{真实世界部署建议。}基于原型经验，我们提炼四条部署原则：\textbf{(a)}~安全策略应支持热加载更新，而非静态配置（规则语法见附录~\ref{sec:appendix-dsl}）；\textbf{(b)}~DCBF 探针维度应匹配部署能力层——白盒/半白盒选用多维探针，黑盒 API 降级为段落级后置过滤（与 \S\ref{sec:physics-regimes} 分层一致）；\textbf{(c)}~置信度加权投票阈值应按业务安全等级校准——高安全场景保持 Deny-优先，高可用场景可放宽 Revise 阈值以减少误拒；\textbf{(d)}~会话交互预算 $N_{\max}$ 与 Token 限额 $T$ 应根据攻击者查询成本独立配置。

\textbf{有限精度 Transformer 与图灵完备建模之间的 Gap。}本文在 Assumption~1 中将 LLM 建模为图灵完备概率程序，以调用 Rice 定理证明全局安全属性的不可判定性。然而，实际部署的 Transformer 是\textbf{有限自动机}：P\'{e}rez et al.~\cite{perez-turing-2021} 的图灵完备性证明依赖无限精度算术与无界序列长度，而 Weiss et al.~\cite{weiss-rnn-counter-2018} 证明有限精度 RNN 等价于计数器机器的有界子类。对于固定精度、固定上下文窗口的 Transformer，其状态空间是有限的，因而安全属性在原则上是可判定的——但状态空间的规模使得穷举验证在计算上不可行（指数级于参数位数 $\times$ 上下文长度）。因此，图灵完备建模的实际意义在于：它为安全分析提供了一个\textbf{保守上界}——即使实际系统``可判定''，其验证复杂度仍使得``外部解耦验证 + 实例级多项式检查''成为唯一实用路径。换言之，不可判定性论证的价值不在于精确刻画 Transformer 的计算能力，而在于论证\textbf{寻求全局隐式安全保证的方向是无望的}，从而为解耦范式提供理论动机。

\textbf{Break-glass 在代数框架中的定位。}安全算子代数（Theorem~3.3）要求组合算子单调且保安全；\texttt{BreakGlassPolicy} 允许在验证器冲突时覆盖 Deny 为 Allow，形式上构成一条\textbf{非单调路径}。我们将其建模为\textbf{受监督审计例外}：break-glass 仅在审计链完整记录时方可生效，$\sigma_{\mathrm{audit}} \circ \sigma_{\mathrm{break\text{-}glass}}$ 的组合仍收敛至 fail-safe $\bot$（Theorem~3.6），因为审计链的不可篡改性保证了任何非单调覆盖都有人类锚定的终止条件。高安全场景建议关闭 break-glass 或将其频率限制在极低水平。

\textbf{EmptySafeCandidateSet 处理策略。}当候选集中无 Token 通过安全断言时，系统默认硬拒绝并触发优雅降级；可选的动态 $K$ 扩容策略在高可用场景下缓解误拒，但需部署者显式启用。策略设计、权衡分析与配置指南详见附录~\ref{sec:appendix-empty-candidate}。

\textbf{DSL 局限与扩展方向。}当前声明式安全策略 DSL 在嵌套条件逻辑、跨会话时序约束与多语言模式匹配方面存在已知局限（语法与规则示例见附录~\ref{sec:appendix-dsl}）。未来扩展路径包括脚本沙箱层（Lua/WASM）、网关侧自动翻译预处理与社区贡献的多语种模式库。

\textbf{形式化验证展望。}安全内核的确定性自动机结构天然适合模型检测（CTL/LTL 编码 + SPIN/NuSMV），可验证策略无冲突与状态可达安全。当前此为\textbf{未来工作}，主要挑战在于将 DSL 语义精确映射为自动机转移函数，并处理语义分类器引入的概率转移扩展。

\textbf{持续攻防评测闭环。}静态测试集无法覆盖演化中的攻击技术；我们建议通过生产 near-miss 采集、自动化 Red Teaming（PAIR~\cite{pair-jailbreak}、GCG~\cite{gcg} 变种）与定期攻防周期构建\textbf{动态 Jailbreak 测试集}，将新绕过路径持续转化为 DSL 规则更新或验证器再训练。

%-------------------------------------------------------------------------------
\section{结论}
%-------------------------------------------------------------------------------

我们将 LLM 安全从「防御模块拼盘」重构为\textbf{范式}：\textbf{解耦安全理论}，奠基三大基本失效定理及其正面表述——计算/验证不对称下的多轮防御、价值/本体锚定与防御性降级、元认知中的表征异常。核心理论贡献在于严格证明了全局隐式安全的不可判定性、合取依赖下安全的不可组合性，以及解耦安全算子在多项式有界候选集下的实例级验证可行性与前向不变性保证（其中 DCBF 类前向不变性结论条件于代理映射忠实性，见第~\ref{sec:physics}~章与 \S\ref{sec:problem-statement} 边界声明）。架构重建（上下文混淆引擎、带分布评估的混合裁判集成、防御性降级与异构恢复）与受控变异（元认知监控器、联合键示例检索、高密度专家迁移）是该范式的\emph{技术实例化}。我们以第~\ref{sec:physics} 章收束\textbf{边界条件与失效物理学}——前提、有限保证、退化分层与非目标——使全文 claim 可被部署者与审稿人核验。随着 LLM 成为基础设施，在本文显式前提与部署约束下，强动机意义上的必要性论证（\textit{computational necessity}，见第~\ref{sec:physics}~章）支持将``隐式对齐''转向``解耦安全''，并为机器速度场景下的\textbf{可审计、可收敛部署型安全}提供一条可辩护路径；\textbf{它不承诺}对通用自然语言交互的全局安全证明，亦非对``无条件可扩展、可审计且具韧性''的普适保证。

%-------------------------------------------------------------------------------
\appendix

\section{附录：安全策略 DSL 语法与规则示例}
\label{sec:appendix-dsl}

安全内核采用声明式 JSON 策略 DSL 定义安全规则，支持热加载（RFC \S 8）。每条策略文件包含 \texttt{policy\_id}、版本号、规则数组及元数据。规则类型包括：

\begin{itemize}
\item \texttt{regex\_deny}：正则匹配输入/输出文本，命中即触发 Deny 或 Revise。
\item \texttt{keyword\_flag}：关键词列表匹配，用于低误报的粗粒度过滤。
\item \texttt{semantic\_classifier}：语义分类器调用（如 Llama-Guard），按阈值判定有害类别。
\item \texttt{cross\_session\_guard}：跨会话/跨 Agent 的合取能力依赖检测（Theorem~3.1 defense）。
\item \texttt{budget\_guard}：交互预算限制（轮次、Token 数），约束攻击者查询预算 $B$。
\item \texttt{rate\_limit}：速率限制，触发优雅降级队列。
\end{itemize}

\noindent\textbf{规则类型能力矩阵。}下表总结各规则类型的适用场景、局限性及与 Verifier Ensemble 的互补关系：

\begin{table}[h]
\centering
\small
\setlength{\tabcolsep}{3pt}
\begin{tabular}{@{}p{2.8cm}p{3.5cm}p{3.5cm}p{3.2cm}@{}}
\toprule
\textbf{规则类型} & \textbf{适用场景} & \textbf{局限性} & \textbf{互补机制} \\
\midrule
\texttt{regex\_deny} & PII 检测、已知注入模式、格式化数据泄露 & 易被同义词替换、Unicode 变体、拼写变形绕过 & 语义分类器捕捉同义表达 \\
\texttt{keyword\_flag} & 低误报的粗粒度过滤、已知敏感词 & 无法处理隐喻、暗语、跨语言表达 & DCBF 探针检测语义偏移 \\
\texttt{semantic\_classifier} & 有害内容分类、上下文感知判定 & 依赖分类器精度、可能误判边缘案例 & 多裁判投票降低单分类器偏差 \\
\texttt{cross\_session\_guard} & 多 Agent 合取攻击防御 & 需要跨会话状态共享基础设施 & 超图模型提供理论保证 \\
\texttt{budget\_guard} & 交互预算控制、查询限速 & 对单次高效攻击无防护 & 语义探针检测单步违规 \\
\texttt{rate\_limit} & DoS 防护、资源配额 & 合法高频用户可能被误限 & 降级 FSM 提供优雅退化 \\
\bottomrule
\end{tabular}
\caption{DSL 规则类型能力矩阵。每种规则类型都有明确的适用边界，安全内核通过规则 + Verifier Ensemble + DCBF 探针的三层互补实现纵深防御。}
\label{tab:dsl-capability-matrix}
\end{table}

\noindent\textbf{示例 1：PII 过滤规则（信用卡号检测）。}

\begin{verbatim}
{
  "rule_id": "deny-credit-card",
  "type": "regex_deny",
  "pattern": "\\b(?:4[0-9]{12}(?:[0-9]{3})?|
               5[1-5][0-9]{14})\\b",
  "scope": "output",
  "action": "deny",
  "severity": "critical"
}
\end{verbatim}

\noindent\textbf{示例 2：提示注入防御规则（忽略指令模式）。}

\begin{verbatim}
{
  "rule_id": "deny-ignore-instructions",
  "type": "regex_deny",
  "pattern": "(?i)(ignore|disregard)\\s+
               (previous|prior)\\s+instructions",
  "scope": "input",
  "action": "deny",
  "severity": "critical"
}
\end{verbatim}

\noindent\textbf{示例 3：多 Agent 合取能力围栏。}

\begin{verbatim}
{
  "rule_id": "deny-cross-agent-credential-relay",
  "type": "cross_session_guard",
  "pattern": "conjunctive_dependency",
  "max_shared_capabilities": 1,
  "scope": "inter_agent",
  "action": "deny",
  "severity": "critical"
}
\end{verbatim}

完整策略示例见仓库 \texttt{rfc\_contracts/example\_policies/} 目录（共 6 个策略文件覆盖 PII 过滤、提示注入防御、有害内容分类、多 Agent 能力围栏、交互预算控制和阶梯式加严策略）。

\noindent\textbf{示例 4：阶梯式加严组合策略（\texttt{escalating\_sensitivity.policy.json}）。}该策略展示 DSL 的\textbf{规则组合能力}——四条规则协同实现"多轮机密请求阶梯式加严"防御：\textbf{(a)} \texttt{cross\_session\_guard} 检测同一会话中对多类凭证（密码、API 密钥、证书等）的合取请求模式，触发即时 Deny；\textbf{(b)} \texttt{rate\_limit} 对敏感主题查询实施 60 秒窗口内最多 5 次的速率限制，超限时返回 Revise（降级提示）；\textbf{(c)} \texttt{budget\_guard} 对每会话敏感交互设置硬上限 10 次，超限后所有敏感请求被 Deny；\textbf{(d)} \texttt{keyword\_flag} 检测社会工程学常用话术（如"经理已授权"、"紧急覆写"等），返回 Revise 提示。四条规则按优先级叠加：速率限制提供渐进预警，预算限制提供硬约束，跨会话守卫提供语义层保护，关键词标记捕捉社工攻击。此组合策略的每条规则独立可审计（Evidence Chain 记录各规则的 \texttt{audit\_tag}），展示了 DSL 声明式组合在复杂安全场景中的表达能力。

\section{附录：多探针屏障值模拟可视化}
\label{sec:appendix-probe-viz}

下表展示一次 DAN（"Do Anything Now"）越狱攻击场景下，三类 DCBF 探针屏障值 $h_1, h_2, h_3$ 随生成步 $t$ 的变化模拟数据。模拟参数：$\theta_{\mathrm{safe}}=0.3$，$\theta_{\mathrm{smooth}}=0.5$，$\theta_{\mathrm{entropy}}=0.3$；中断条件为任一 $h_i < 0$。

\begin{table}[h]
\centering
\small
\begin{tabular}{@{}r r r r l@{}}
\toprule
\textbf{步 $t$} & $h_1$ (语义) & $h_2$ (轨迹) & $h_3$ (熵) & \textbf{状态} \\
\midrule
1 & 0.22 & 0.41 & 0.26 & safe \\
2 & 0.19 & 0.38 & 0.24 & safe \\
3 & 0.15 & 0.33 & 0.21 & safe \\
4 & 0.10 & 0.28 & 0.18 & safe \\
5 & 0.05 & 0.19 & 0.11 & safe \\
6 & $-0.03$ & 0.08 & 0.04 & \textbf{INTERRUPT ($h_1<0$)} \\
\bottomrule
\end{tabular}
\caption{DAN 越狱场景多探针屏障值模拟。攻击者在 $t=1$ 注入角色征用提示后，语义探针 $h_1$ 持续下降，在 $t=6$ 首先越过零界面触发中断。轨迹平滑性探针 $h_2$ 和熵探针 $h_3$ 同步下降但略滞后。三探针 AND 共识下，$h_1$ 的突破即足以拦截攻击。}
\label{tab:probe-sim-dan}
\end{table}

\noindent\textbf{解读。}该模拟数据展示了三探针的互补特性：语义探针 $h_1$ 对语义偏离最为敏感，率先检测到越狱企图；轨迹平滑性探针 $h_2$ 捕捉嵌入跳跃（如角色切换导致的嵌入突变）；熵探针 $h_3$ 检测输出概率分布的突然变化（如从正常回答切换到指令复述）。对于渐进式攻击（如 Crescendo），单步 $h_i$ 可能缓慢逼近零但不越界，此时 CUSUM 滑动窗口监控器（\S 附录~\ref{sec:appendix-probe-viz}）可通过累积统计量捕捉整体下行趋势。

\section{附录：裁判集成置信度来源与校准}
\label{sec:appendix-confidence}

各类验证器的 confidence 值来源不同：\textbf{(a)} LLM 型验证器（如 Llama-Guard）使用其输出 softmax 概率中安全/不安全类别的最大概率作为 $c_j$。\textbf{(b)} 规则型验证器（如 regex\_deny、keyword\_flag）因判定为确定性匹配，固定 $c_j = 1.0$。\textbf{(c)} 外部分类器（如 StrongREJECT 评分器）使用其评分归一化至 $[0,1]$ 区间。为缓解 LLM 验证器的\textbf{过度自信}问题，可选启用温度缩放后处理：$c_j' = \sigma(\log(c_j) / T)$，其中 $T > 1$ 为校准温度参数，在验证集上通过最小化预期校准误差（ECE）获得。关键不变量：Deny-优先原则确保即使单个验证器给出高 confidence 的 Allow 判决，只要存在\emph{任意} Deny 票，最终裁决仍为 Deny。该不变量使加权投票的安全保证\textbf{单调}——增加验证器数量或调整 confidence 仅可能使系统更严格，不会意外放松。

\textbf{典型验证器性能参考。}Llama-Guard 系列在 JBB-Behaviors 上的 F1 约 0.82--0.89，良性误拒率（FPR）$<$5\%~\cite{llama-guard}；StrongREJECT 评分器在压制空越狱后精度约 0.91~\cite{strongreject}；结合多裁判投票后系统级 FPR 可降至 $<$2\%（详见评估 \S\ref{sec:evaluation}）。

\section{附录：EmptySafeCandidateSet 处理策略}
\label{sec:appendix-empty-candidate}

当 Axiom Hive 在 Top-$K$ 候选集中找不到任何通过安全断言的 Token（\texttt{EmptySafeCandidateSet}）时，默认策略为\textbf{硬拒绝}——返回空输出并触发优雅降级 FSM（DegradedShutdown）。这是\textbf{安全保守}的设计选择：宁可误拒（保护用户），不冒险放行未经验证的 Token。

\textbf{动态 $K$ 扩容策略。}为缓解过度保守对可用性的损害，代码库实现了可选的 \texttt{DynamicKExpansionPolicy}：当 \texttt{EmptySafeCandidateSet} 触发时，将 $K$ 翻倍（最大至 $K_{\max}=512$）并重新执行投影，至多重试 2 次。该策略默认\textbf{关闭}，需由部署者根据业务场景显式启用。

\textbf{权衡分析。}硬拒绝策略适合高安全场景（金融、医疗），因为误放行的代价远高于误拒绝；动态扩容适合高可用场景（客服、创作辅助），因为用户体验对空回复高度敏感。两种策略的选择是部署时的\textbf{配置决策}，而非架构限制——安全内核的可审计性保证在两种模式下均成立（扩容过程的每次重试均记录在 Evidence Chain 中）。

\section{附录：主流 LLM 安全机制横向对比}
\label{sec:appendix-comparison}

\begin{table*}[t]
\centering
\small
\setlength{\tabcolsep}{3pt}
\begin{tabular}{@{}p{2.8cm}p{2.4cm}p{2.2cm}p{2.2cm}p{2.0cm}p{2.0cm}@{}}
\toprule
\textbf{方案} & \textbf{拦截机制} & \textbf{介入粒度} & \textbf{规则灵活性} & \textbf{可审计性} & \textbf{热更新} \\
\midrule
OpenAI Moderation API & 分类器后置过滤 & 段落级 & 固定类别、不可定制 & 仅返回聚合分数 & 不透明 \\
RLHF / RLAIF & 训练时对齐 & 模型参数级 & 需重新训练 & 无运行时审计 & 需重训 \\
Constitutional AI (Anthropic) & 原则引导自我纠正 & 模型生成级 & 宪法原则可编辑 & 原则链可追溯 & 需重训 \\
Google 内容过滤 & 分类器 + 关键词 & 段落级 & 有限可配置 & 部分分类日志 & 有限热更新 \\
Llama-Guard & 安全分类器后置 & 段落级 & 自定义分类法 & 分类标签可追溯 & 模型更新 \\
\textbf{解耦安全内核}（本文） & \textbf{多层纵深}（DSL + Judge Ensemble + DCBF + 超图） & \textbf{Token 级}（白盒）/ 段落级（黑盒） & \textbf{DSL 热加载}、规则任意组合 & \textbf{Evidence Chain} 全链路审计 & \textbf{分钟级} DSL 热更新 \\
\bottomrule
\end{tabular}
\caption{主流 LLM 安全机制横向对比。解耦安全内核在介入粒度、规则灵活性、可审计性和热更新能力四个维度上具有结构性优势，但以额外架构复杂度和部署成本为代价。}
\label{tab:safety-comparison}
\end{table*}

\section{附录：评估矩阵与构建配方}
\label{sec:appendix-eval}

表~\ref{tab:eval-matrix}汇总威胁模型轨道、典型攻击/数据集与主指标对应关系；构建 PDF 时须将所用模型、检查点、依赖版本与随机种子写入制品元数据（如 \texttt{sbom.json} 或实验日志头），以满足可复现性审查。

\begin{table*}[t]
\centering
\small
\setlength{\tabcolsep}{4pt}
\begin{tabular}{@{}p{2.6cm}p{2.4cm}p{3.4cm}p{5.8cm}@{}}
\toprule
\textbf{轨道} & \textbf{攻击/压力源} & \textbf{数据与预算} & \textbf{主报告指标} \\
\midrule
Track A & 多轮抽取、HPM、越狱提示、诱饵 DoS & 标准有害集/多轮基准；固定 $N_{\max}$、Token 预算、速率限制 & ASR、RSR、\textbf{良性任务成功率}、残余抽取 F1、\textbf{良性 SLA}（\S\ref{sec:eval-utility}） \\
Track B & GCG、表征对齐 PBU & 同型 victim；声明梯度步数、对抗串长度、可查询隐藏状态与否 & 组件上界 ASR、攻击成本 $\hat{C}$；\textbf{不与} Track A 混排 \\
\midrule
\multicolumn{4}{@{}p{0.96\textwidth}@{}}{\textbf{复现清单（应在日志/附录中逐项填写）：}基座模型与 revision；Tokenizer/聊天模板；各裁判模型 ID 与版本；上下文混淆与诱饵参数哈希；联合键阈值与良性验证集版本；解码温度与 top-$p$；随机种子列表；Track A/B 与披露等级 L0--L2。} \\
\bottomrule
\end{tabular}
\caption{评估矩阵：\textbf{Track A/B} 分轨、典型攻击面与主指标。良性 SLA 与效用切片见 \S\ref{sec:eval-utility}。}
\label{tab:eval-matrix}
\end{table*}

\begin{table}[t]
\centering
\small
\begin{tabular}{@{}ll@{}}
\toprule
\textbf{提示文件（仓库路径）} & \textbf{SHA256} \\
\midrule
\texttt{.../data/harmful\_prompts\_trackA\_en.txt} & \texttt{55f1281e97d7fd839d6d707ae66a6857ae54dd7cfded994d76fe13dd1bfcc568} \\
\texttt{.../data/benign\_prompts\_en.txt} & \texttt{39c48b0022a89514542d2e7f141ad54ccbecf30817cdb34efc03e1ff3f048917} \\
\texttt{.../data/hpm\_proxy\_prompts\_en.txt} & \texttt{67b5ac7b36cf8cfcad64319715057268abd38be2a71d7049ac16418fe856a51d} \\
\bottomrule
\end{tabular}
\caption{Track~A 固定提示表（\texttt{Decoupled-LLM-Gateway/experiments/data/}，完整路径见 \texttt{Data\_Mapping\_TrackA\_to\_Paper.md}）。修改后须重算哈希。}
\label{tab:prompt-sha}
\end{table}

\section{附录：评估制品与复现细节}
\label{sec:appendix-eval-artifact}

本节承载 \S\ref{sec:eval-artifact} 主文中下沉的工程细节，便于 AE/复现审查，同时避免主文被实现说明淹没。

\noindent 配套仓库中的 \texttt{Decoupled-LLM-Gateway} 为\textbf{参考实现}：同步网关（上下文混淆、诱饵注入、策略驱动的模板降级）与异步 Redis 流、策略 Hash 刷新及 Python Worker（当前为\textbf{诱饵泄露}裁判 MVP，可扩展为多模型集成）。\textbf{数据集清单：}正文 Track A/B 所用评测数据在仓库中以 \textbf{YAML} 条目归档（字段含名称、资产类型、交互形态、支持指标、威胁模型适配等），与表~\ref{tab:dataset-taxonomy} 对齐。评测脚本 \texttt{experiments/run\_paper\_benchmark.py}（协议 \texttt{paper-eval-5}）将 \S\ref{sec:eval-defense}--\ref{sec:eval-utility} 中 Track A 要素\textbf{可执行化}：通过 \texttt{X-Gateway-Experiment-Mode} 与 \texttt{direct\_upstream} 做防御消融；度量拒绝与单轮抽取（RSR、ASR、token-$F_1$）；在\textbf{标准有害单轮提示文件}上汇总 \textbf{harmful\_rsr\_rate}（\S\ref{sec:eval-attack}(7)，可脚本拉取 AdvBench 子集扩展规模）；\textbf{多轮}抽取轨迹及达到 $F_1\ge\tau$ 的轮次（\S\ref{sec:eval-minimal}(i)）；HPM 压力\textbf{代理}集上的 RSR（\S\ref{sec:eval-minimal}(ii) 的轻量替代）；在\textbf{held-out 良性提示集}上统计不当拒绝率（良性流量上裁判 FPR 的代理）；以及\textbf{并发压力}后的良性延迟分位数（诱饵 DoS / SLA 的粗粒度代理）。

\noindent \texttt{manifest} 记录 \texttt{judge\_mode}、\texttt{smooth\_llm\_samples\_k} 等；HTTP 裁判可由 \texttt{experiments/judge\_service/server.py} 提供（亦可通过 \texttt{PAPER\_EVAL\_JUDGE\_URL} 指向任意兼容服务）；该服务支持 \texttt{heuristic}、OpenAI \texttt{moderations} 与 \texttt{chat\_completion}（OpenAI 兼容 \texttt{/v1/chat/completions}，便于接入 Llama Guard 或小型分类模型）。\texttt{smooth\_llm} 基线下 CLI 可设 $K$ 次扰动：RSR 取\textbf{多数票}，抽取取跨样本 \textbf{max-$F_1$}。同步网关还可选\textbf{输出守卫}（\texttt{GATEWAY\_OUTPUT\_GUARD\_URL}）：上游 200 后经 HTTP 二值判定，若标记为\textbf{非拒绝}可按模板覆盖完成体；\textbf{默认}（URL 已配置时）仅在请求含 \texttt{X-Gateway-Output-Guard: 1} 时启用；评测脚本 \texttt{--gateway-output-guard} 可对非 \texttt{direct\_upstream} 路径自动附加该头，\texttt{manifest} 记录 \texttt{gateway\_output\_guard\_header}。

\noindent 将 \texttt{GATEWAY\_UPSTREAM} 指向真实 OpenAI 兼容 API 可生成\textbf{正文主表用的实证数据}；在 \texttt{Decoupled-LLM-Gateway/} 下先执行 \texttt{. ./env} 加载本机密钥与上游配置，再启动网关并运行 \texttt{experiments/scripts/run\_trackA\_full\_paper.sh}。\texttt{echo-llm} 与 \texttt{experiments/scripts/run\_trackA\_full\_echo\_ci.sh} 可在相同场景矩阵下做低成本回归与对照实验。\textbf{Track B}（GCG、表征对齐 PBU）\textbf{不在}该制品范围内（\S\ref{sec:eval-tracks}）。可用 \texttt{python3 experiments/scripts/export\_trackA\_table\_latex.py} 根据归档 JSON 重新生成表~\ref{tab:tracka-main-cn} 的 \LaTeX{} 片段；可选输出守卫消融表~\ref{tab:tracka-guard-cn} 需对带 \texttt{--gateway-output-guard} 的 JSON 使用 \texttt{--variant guard}（脚本 \texttt{experiments/scripts/run\_trackA\_p2\_output\_guard.sh}）。实验建议与表~\ref{tab:tracka-main-cn}、表~\ref{tab:tracka-guard-cn}、表~\ref{tab:eval-matrix} 一并报告模型 ID、随机种子与原始 JSON。

%-------------------------------------------------------------------------------
\section*{可用性}
%-------------------------------------------------------------------------------
制品与扩展证明将在经机构与出口管制审查后提供；构建产物（含 \texttt{.tex}、\texttt{.bib}、生成 PDF 与复现脚本）默认置于 \texttt{paper/} 目录以便版本对齐。参考网关与 \texttt{paper-eval-5} 评测脚本（含 \texttt{judge\_service}）位于 \texttt{Decoupled-LLM-Gateway/}（\S\ref{sec:eval-artifact}）。

%-------------------------------------------------------------------------------
\section*{致谢}
%-------------------------------------------------------------------------------
感谢产业与标准界同仁在安全架构与数据治理方面的反馈。

%-------------------------------------------------------------------------------
\bibliographystyle{plain}
\bibliography{ArchitecturalReconstruction_ControlledMutation}

\end{CJK}
\end{document}
